% !TeX program = xelatex
% BẢN THẢO V1 — nội dung được tổng hợp từ mã nguồn và báo cáo thí nghiệm trong repo.
% TODO trước khi nộp: thay thông tin trường, tác giả, giảng viên hướng dẫn và chuẩn trích dẫn.

\documentclass[12pt,a4paper,oneside]{report}

\usepackage{iftex}
\ifPDFTeX
  \usepackage[utf8]{inputenc}
  \usepackage[T5]{fontenc}
  \usepackage[vietnamese]{babel}
\else
  \usepackage{fontspec}
  \setmainfont{Times New Roman}
  \setsansfont{Arial}
  \setmonofont{Consolas}
\fi

\usepackage[a4paper,top=2.5cm,bottom=2.5cm,left=3.2cm,right=2.3cm]{geometry}
\usepackage{amsmath,amssymb,mathtools}
\usepackage{booktabs,longtable,tabularx,array,multirow}
\usepackage{graphicx,float}
\usepackage{xcolor}
\usepackage{enumitem}
\usepackage{listings}
\usepackage{caption}
\usepackage{fancyhdr}
\usepackage{microtype}
\usepackage{setspace}
\usepackage{url}
\usepackage[hidelinks,unicode]{hyperref}

\definecolor{lightgray}{RGB}{245,245,245}
\definecolor{darkred}{RGB}{150,35,35}

\hypersetup{
  pdftitle={Mô hình lai dựa trên khuôn mẫu và học sâu cho dự đoán cấu trúc ba chiều RNA},
  pdfauthor={TODO: Họ tên học viên},
  pdfsubject={Bản thảo luận văn thạc sĩ V1},
  pdfkeywords={RNA 3D, TBM, DRfold2, Geometry v2, temporal safety, TM-score}
}

\setstretch{1.35}
\setlength{\parindent}{1.2cm}
\setlength{\parskip}{0.25em}
\setlist[itemize]{topsep=0.3em,itemsep=0.2em}
\setlist[enumerate]{topsep=0.3em,itemsep=0.2em}
\renewcommand{\arraystretch}{1.18}

\pagestyle{fancy}
\fancyhf{}
\fancyhead[L]{\small Bản thảo luận văn V1}
\fancyhead[R]{\small Mô hình lai dự đoán cấu trúc ba chiều RNA}
\fancyfoot[C]{\thepage}
\setlength{\headheight}{15pt}

\lstset{
  basicstyle=\ttfamily\small,
  backgroundcolor=\color{lightgray},
  frame=single,
  breaklines=true,
  columns=fullflexible,
  keepspaces=true,
  showstringspaces=false
}

\newcommand{\cprime}{C1$^{\prime}$}
\newcommand{\angstrom}{\text{\AA}}
\newcommand{\TODO}[1]{\textcolor{darkred}{\textbf{TODO: #1}}}

\begin{document}

% -----------------------------------------------------------------------------
% TITLE PAGE
% -----------------------------------------------------------------------------
\begin{titlepage}
\begin{center}
  {\large \TODO{UNIVERSITY OF SCIENCE AND TECHNOLOGY OF HANOI}}\\[0.3cm]
  {\large \TODO{ICT LAB}}\\[2.0cm]

  {\Large\bfseries \TODO{ĐỖ THÀNH ĐẠT}}\\[1.5cm]

  {\LARGE\bfseries
  MÔ HÌNH LAI DỰA TRÊN KHUÔN MẪU\\
  VÀ HỌC SÂU CHO DỰ ĐOÁN CẤU TRÚC\\
  BA CHIỀU RNA VỚI KIỂM SOÁT RÒ RỈ\\
  THỜI GIAN VÀ TINH CHỈNH HÌNH HỌC
  }\\[1.0cm]

  {\large\itshape
  A Hybrid Template-Based and Deep-Learning Framework for RNA Three-Dimensional
  Structure Prediction with Temporal Leakage Control and Conservative Geometric Refinement
  }\\[1.5cm]

  {\large LUẬN VĂN THẠC SĨ}\\[0.4cm]
  {\large \TODO{DATA SCIENCE}}\\[1.6cm]

  \begin{tabular}{ll}
  Giảng viên hướng dẫn: & \TODO{NGUYỄN MINH HƯƠNG + NGUYỄN CẨM LINH}\\
  Mã số học viên:       & \TODO{2440059}
  \end{tabular}

  \vfill
  {\large \TODO{HÀ NỘI}, 2026}
\end{center}
\end{titlepage}

\pagenumbering{roman}

% -----------------------------------------------------------------------------
% FRONT MATTER
% -----------------------------------------------------------------------------
\chapter*{Lời cam đoan}
\addcontentsline{toc}{chapter}{Lời cam đoan}
\TODO{Bổ sung lời cam đoan theo mẫu chính thức của trường.}

\chapter*{Lời cảm ơn}
\addcontentsline{toc}{chapter}{Lời cảm ơn}
\TODO{Bổ sung lời cảm ơn giảng viên hướng dẫn, đơn vị đào tạo, đồng nghiệp và gia đình.}

\chapter*{Tóm tắt}
\addcontentsline{toc}{chapter}{Tóm tắt}

Dự đoán cấu trúc ba chiều RNA từ trình tự vẫn khó do dữ liệu thực nghiệm ít, RNA linh
động và một trình tự có thể có nhiều cấu dạng. Luận văn nghiên cứu một quy trình dự
đoán ở mức nguyên tử \cprime{} gồm mô hình hóa dựa trên khuôn mẫu (TBM), ứng viên từ
mô hình học sâu DRfold2 và một phép tinh chỉnh bảo thủ tên Geometry v2. Mục tiêu chính
là sinh năm giả thuyết có TM-score cao; câu hỏi phương pháp là liệu sửa hình học cục bộ
có tăng độ đúng cục bộ mà không làm hỏng nếp gấp tổng thể hay không.

Mọi khuôn mẫu được lọc theo nguyên tắc phát hành trước đối tượng dự đoán và loại PDB của
chính đối tượng. Phân tích TBM cho thấy MMseqs2 chỉ tìm được khuôn ở 8/20 RNA hiệu chỉnh,
trong khi tìm kiếm composite tạo ứng viên cho 20/20. Xếp hạng bằng identity nhân coverage
tăng best-of-five TM trung bình 0.01566 so với identity đơn lẻ, còn thành phần
completeness và luật distinct-PDB không tăng TM trên cohort này. Công thức composite
$0.4G+0.3L+0.2F+0.1K_3$ ổn định trước thay đổi trọng số nhưng không tốt hơn rõ ràng
so với chỉ dùng căn chỉnh toàn cục; do đó luận văn xem nó là heuristic tăng recall,
không khẳng định tối ưu. Bỏ lọc thời gian làm điểm tăng lạc quan khoảng 0.02696.

Nhánh TBM đạt 0.60084 công khai và 0.60175 riêng tư trên Kaggle. Quy trình ba TBM, hai
DRfold2 và Geometry v2 đạt 0.62809/0.61390. Trên 20 RNA mới nhất, TBM oracle đạt
0.565183, DRfold2 oracle 0.502627, nhưng hợp hai nguồn đạt 0.588304: DRfold2 thắng TBM
ở 8/20 RNA và bổ sung nếp gấp thay vì thay thế TBM. Với cùng ba vị trí dự đoán, thay một
TBM bằng một DRfold2 tăng TM trung bình 0.016012 nhưng khoảng tin cậy còn chứa 0.

Đánh giá Geometry v2 dùng giao thức thời gian và họ trình tự 60/20/20: 60 RNA ước lượng
phân bố hình học, 20 RNA chọn tham số, rồi cấu hình được khóa trước khi mở 20 RNA mới
nhất. Trên tập xác nhận, \cprime{}-lDDT tăng từ 0.659836 lên 0.666496
($\Delta=+0.006660$, khoảng tin cậy 95\% $[+0.000933,+0.013093]$); RMSD cửa sổ 9 và
15 giảm lần lượt 0.040386 và 0.017257~\angstrom. TM-score thay đổi từ 0.588304 xuống
0.587197 ($\Delta=-0.001107$, khoảng tin cậy chứa 0), nằm trong biên bảo toàn 0.005.
Đối chứng chỉ giữ nguồn và khoảng cách backbone tăng lDDT tương đương, nhưng Geometry
v2 cải thiện mạnh hơn RMSD cửa sổ 9, giảm va chạm, NLL góc và NLL góc xoắn, đồng thời
không làm tăng góc gập sắc. Lợi ích lDDT tập trung ở ứng viên ban đầu yếu; nhóm chất
lượng cao không có bằng chứng cải thiện.

Kết luận chính là: tìm kiếm composite mở rộng độ phủ khuôn; TBM và DRfold2 có tính bổ
sung; Geometry v2 tạo cải thiện cục bộ nhỏ nhưng có ý nghĩa thống kê trong khi bảo toàn
nếp gấp tổng thể theo biên định trước. GeoFuse được giữ như thí nghiệm khám phá trong
phụ lục và không được xem là đóng góp lõi vì chưa tạo cải thiện khái quát.

\noindent\textbf{Từ khóa:} cấu trúc ba chiều RNA, mô hình hóa dựa trên khuôn mẫu,
DRfold2, TM-score, \cprime{}-lDDT, an toàn thời gian, tinh chỉnh hình học.

\tableofcontents
\listoftables
\listoffigures
\clearpage
\pagenumbering{arabic}

% -----------------------------------------------------------------------------
% CHAPTER 1
% -----------------------------------------------------------------------------
\chapter{Giới thiệu}

\section{Bối cảnh và động lực nghiên cứu}

RNA tham gia nhiều quá trình sinh học như điều hòa biểu hiện gene, xúc tác, nhận diện
phân tử và hình thành các phức hợp ribonucleoprotein. Những chức năng này không chỉ phụ
thuộc vào trình tự A, C, G và U mà còn phụ thuộc vào cách chuỗi RNA tổ chức thành cấu
trúc bậc hai và cấu trúc ba chiều. Cấu trúc bậc hai mô tả chủ yếu các cặp base nitơ và
các mô-típ như thân xoắn hoặc vòng; cấu trúc bậc ba xác định cách các mô-típ đó sắp xếp,
tiếp xúc và xoắn
trong không gian. Vì vậy, khả năng suy ra cấu trúc ba chiều từ trình tự có ý nghĩa đối
với nghiên cứu chức năng RNA và thiết kế các phân tử tác động lên RNA.

Cấu trúc RNA có thể được xác định bằng tinh thể học tia X, cộng hưởng từ hạt nhân hoặc
kính hiển vi điện tử nghiệm lạnh. Tuy nhiên, các quy trình này thường đòi hỏi nhiều thời
gian, chi phí và không phải RNA nào cũng phù hợp để đo. So với protein, số lượng cấu trúc
RNA thực nghiệm còn hạn chế; hơn nữa RNA có độ linh động cao, nhiều tương tác không chuẩn
Watson--Crick và các motif bậc ba phụ thuộc mạnh vào môi trường. Những đặc điểm đó khiến
dự đoán RNA 3D vẫn là một bài toán chưa được giải quyết hoàn toàn
\cite{rna_puzzles,rna_progress_review,state_of_rnart}.

Các phương pháp tính toán hiện nay có thể được nhìn theo ba nhóm chính. Phương pháp dựa
trên vật lý tìm cấu hình có năng lượng thuận lợi nhưng phải đối mặt với không gian tìm
kiếm rất lớn. Phương pháp dựa trên tri thức hoặc khuôn mẫu tái sử dụng motif và cấu trúc
đã biết; nhóm này thường chính xác khi có khuôn mẫu phù hợp nhưng suy giảm khi gặp nếp
gấp mới. Phương pháp học sâu học mối liên hệ giữa trình tự và hình học từ dữ liệu; chúng
có thể tạo dự đoán không phụ thuộc trực tiếp vào một khuôn mẫu gần, nhưng bị giới hạn
bởi lượng dữ liệu RNA, độ lệch phân bố và khó khăn trong ước lượng độ chính xác của từng
dự đoán \cite{drfold,rhofoldplus,drfold2,ares}.

Như vậy, không có một nguồn dự đoán đơn lẻ nào luôn chiếm ưu thế. Khuôn mẫu thường mạnh
khi có họ hàng gần; mô hình học sâu có thể bổ sung nếp gấp khi khuôn yếu. Luận văn trước
hết hỏi từng nguồn làm tăng độ phủ nếp gấp ở đâu, sau đó tập trung vào một câu hỏi hẹp
hơn: khi nếp gấp ban đầu đã có, liệu có thể sửa quan hệ \cprime{} cục bộ mà không làm
dịch chuyển bố cục toàn phân tử hay không?

Động lực trung tâm vì vậy là xây dựng một pipeline RNA 3D an toàn theo thời gian, phân
tích định lượng thành phần TBM và tính bổ sung TBM--DRfold2, rồi kiểm chứng Geometry v2
bằng ablation và tập xác nhận mới nhất. Luận văn không lấy ước lượng chất lượng cục bộ
hay hợp nhất tọa độ làm contribution lõi; GeoFuse chỉ còn là thí nghiệm khám phá ở phụ
lục vì chưa cho thấy cải thiện khái quát.

\section{Bối cảnh bài toán thực nghiệm}

Cuộc thi Stanford RNA 3D Folding \cite{kaggle_competition} được sử dụng như bối cảnh triển khai và đánh giá bên
ngoài. Với mỗi RNA, hệ thống phải xuất năm cấu trúc, mỗi nucleotide được biểu diễn bằng
tọa độ nguyên tử \cprime{}. Chất lượng được đánh giá bằng TM-score tốt nhất trong năm
dự đoán. Thiết lập này phản ánh hai đặc điểm thực tế: dự đoán cấu trúc có tính bất định,
và khả năng sinh một tập giả thuyết đa dạng có thể quan trọng không kém việc tối ưu một
dự đoán duy nhất.

Các mã nguồn dạng sổ tay tính toán của nhóm đứng đầu cuộc thi cho thấy mô hình hóa dựa trên khuôn mẫu và mô hình học
sâu có thể kết hợp hiệu quả ở cấp độ hệ thống. Tuy nhiên, một giải pháp thi đấu chủ yếu
tối ưu điểm số và thời gian chạy, trong khi một luận văn cần tách rõ dữ liệu nào được
sử dụng, thành phần nào tạo ra cải thiện, kết quả có tổng quát hay không và giới hạn
sinh học của biểu diễn \cprime{}. Do đó luận văn không xem bảng xếp hạng là bằng chứng
duy nhất, mà xây dựng các tập phát triển và xác nhận riêng với kiểm soát thời gian, họ
trình tự và thống kê theo đối tượng RNA.

\section{Phát biểu bài toán}

Cho một RNA cần dự đoán có trình tự

\begin{equation}
S=(s_1,s_2,\ldots,s_L),\qquad s_i\in\{A,C,G,U\},
\end{equation}

và mốc thời gian $t_{cutoff}$, hệ thống phải sinh năm cấu trúc ứng viên

\begin{equation}
\mathcal{X}=\{X^{(1)},\ldots,X^{(5)}\},\qquad
X^{(k)}\in\mathbb{R}^{L\times 3},
\end{equation}

trong đó hàng thứ $i$ của $X^{(k)}$ là tọa độ $(x_i,y_i,z_i)$ của nguyên tử
\cprime{} đại diện cho nucleotide thứ $i$. Cấu trúc thực nghiệm của RNA không được sử
dụng khi suy luận. Mọi khuôn mẫu phải có ngày phát hành nhỏ hơn $t_{cutoff}$.

Mục tiêu theo thiết lập của cuộc thi là tối đa hóa

\begin{equation}
\frac{1}{N}\sum_{n=1}^{N}
\max_{k\in\{1,\ldots,5\}}
\mathrm{TM}\!\left(X_n^{(k)},Y_n\right),
\end{equation}

với $Y_n$ là cấu trúc thực nghiệm chỉ được sử dụng khi đánh giá.

Ngoài mục tiêu toàn cục trên, luận văn xét thêm độ chính xác cục bộ. Điều này xuất phát
từ nhận xét rằng một dự đoán có nếp gấp đúng vẫn có thể chứa khoảng cách hoặc góc cục
bộ chưa chính xác, trong khi một cấu trúc có hình học trơn tru chưa chắc thuộc đúng nếp
gấp. Hai loại chất lượng này cần được đo bằng các chỉ số bổ sung thay vì quy về một số
duy nhất.

\section{Khoảng trống nghiên cứu và giả thuyết}

Ba khoảng trống được xét. Thứ nhất, điểm toàn hệ thống không cho biết thành phần nào của
TBM tạo recall, precision hay đa dạng. Thứ hai, mức tăng của pipeline lai có thể bị
nhiễu giữa số lượng candidate, nguồn DRfold2 và Geometry. Thứ ba, các chỉ số hình học
nội bộ không đủ chứng minh cấu trúc gần native hơn vì chúng cũng là mục tiêu được tối ưu.

Ba giả thuyết tương ứng là: tìm kiếm composite tăng độ phủ khi MMseqs2 không có hit;
DRfold2 bổ sung chứ không thay thế ứng viên TBM; và Geometry v2 có thể cải thiện các chỉ
số native độc lập ở quy mô cục bộ trong khi giữ TM-score trong biên suy giảm 0.005.

\section{Câu hỏi nghiên cứu}

Luận văn trả lời ba nhóm câu hỏi:

\begin{description}[leftmargin=2.2cm,style=nextline]
\item[RQ1.] Việc kết hợp truy hồi nhanh với tìm kiếm tương đồng mở rộng có cải thiện
khả năng tìm khuôn mẫu trong điều kiện an toàn theo thời gian hay không?

\item[RQ2.] DRfold2 có bổ sung nếp gấp mà TBM bỏ lỡ hay không; hiệu ứng nguồn còn tồn
tại khi giữ cố định số lượng cấu trúc dự đoán hay không?

\item[RQ3a.] Geometry v2 có cải thiện độ chính xác cục bộ theo \cprime{}-lDDT và RMSD
cửa sổ, đồng thời bảo toàn TM-score hay không?
\item[RQ3b.] Thành phần source, backbone, angle, torsion, clash/Rg, kink guard và adaptive
strength đóng vai trò gì so với một đối chứng source--backbone đơn giản?
\item[RQ3c.] Hiệu ứng có khái quát sang 20 RNA mới nhất tách theo thời gian/họ và thay
đổi theo nguồn, chất lượng ban đầu, độ dài hay độ phủ như thế nào?
\end{description}

\section{Mục tiêu}

\subsection{Mục tiêu tổng quát}

Xây dựng và đánh giá một pipeline lai dự đoán RNA 3D mức \cprime{} có thể tái lập và an
toàn theo thời gian; đồng thời kiểm chứng liệu tinh chỉnh hình học bảo thủ có cải thiện
độ đúng cục bộ mà không phá nếp gấp toàn cục hay không.

\subsection{Mục tiêu cụ thể}

\begin{enumerate}
  \item Xây dựng thư viện khuôn mẫu RNA từ PDB/mmCIF và kiểm tra ngày công bố.
  \item Kết hợp MMseqs2 với tìm kiếm tương đồng mở rộng để tăng khả năng truy hồi khuôn mẫu.
  \item Chuẩn hóa đầu ra TBM và DRfold2 thành một biểu diễn ứng viên thống nhất; Boltz
  chỉ được ghi nhận như feasibility case cho RNA dài.
  \item Thiết kế Geometry v2 với các phân bố hình học thực nghiệm và ràng buộc nguồn
  thích nghi theo độ tin cậy.
  \item Phân rã ảnh hưởng của retrieval, reranking, distinct-PDB, gap filling và phân
  bổ TBM--DRfold2 bằng các ablation giữ biến còn lại cố định khi có thể.
  \item Đánh giá Geometry v2 bằng protocol 60 train--20 calibration--20 validation,
  chỉ số native độc lập và bootstrap theo đối tượng RNA.
\end{enumerate}

\section{Đóng góp}

Luận văn có một đóng góp phương pháp và hai đóng góp thực nghiệm chính.

\begin{enumerate}
  \item \textbf{Pipeline TBM an toàn theo thời gian và phân tích white-box.} Luận văn
  kết hợp MMseqs2 với tìm kiếm composite, lọc theo ngày/PDB, căn chỉnh, reranking,
  distinct-PDB và gap filling; mỗi component được gắn với giả thuyết và ablation.

  \item \textbf{Bằng chứng thực nghiệm về tính bổ sung TBM--DRfold2.} Các allocation
  fixed-size tách lợi ích nguồn khỏi lợi thế có thêm prediction slot. Kết quả cho thấy
  DRfold2 yếu hơn TBM trung bình nhưng thắng trên 8/20 RNA và tăng union oracle.

  \item \textbf{Geometry v2 và kiểm chứng confirmatory.} Phương pháp kết hợp source
  restraint, backbone, clash/Rg, angle/torsion empirical NLL, kink guard và adaptive
  strength. Priors chỉ học từ 60 RNA sớm; tham số chọn trên 20 calibration; hiệu quả
  chấm đúng một lần trên 20 RNA mới nhất bằng lDDT/RMSD cửa sổ/TM độc lập.
\end{enumerate}

GeoFuse không còn là contribution: kết quả âm được rút khỏi mạch lập luận chính và chỉ
giữ ở phụ lục như một hướng khám phá đã đóng băng.

\section{Phạm vi}

Luận văn dự đoán đường \cprime{}, không phải cấu trúc RNA toàn nguyên tử. Vì vậy nghiên
cứu có thể kết luận về nếp gấp tổng thể, tọa độ \cprime{}, khả năng bảo toàn khoảng cách
cục bộ và hình học thô; không kết luận trực tiếp về liên kết hydro, xếp chồng base nitơ, cấu
hình đường, hóa học lập thể của phosphate, năng lượng tự do hay độ ổn định trong mô
phỏng động lực học phân tử.

\section{Cấu trúc luận văn}

Chương 2 trình bày nền tảng và tổng quan công trình liên quan. Chương 3 mô tả dữ liệu,
kiểm soát rò rỉ và quy trình đánh giá. Chương 4 trình bày phương pháp lai đề xuất.
Chương 5 mô tả thiết kế thí nghiệm. Chương 6 lần lượt trả lời ba nhóm câu hỏi nghiên cứu.
Chương 7 diễn giải ý nghĩa của các kết quả, giới hạn và quan hệ với công trình trước.
Chương 8 kết luận và đề xuất hướng phát triển.

% -----------------------------------------------------------------------------
% CHAPTER 2
% -----------------------------------------------------------------------------
\chapter{Kiến thức nền và tổng quan công trình liên quan}

\section{Từ trình tự đến cấu trúc ba chiều RNA}

Cấu trúc RNA thường được mô tả theo một hệ phân cấp. Cấu trúc bậc một là trình tự
nucleotide. Cấu trúc bậc hai mô tả các cặp base nitơ và các phần tử như thân xoắn
(stem), vòng kẹp tóc (hairpin), chỗ phình (bulge) hay vòng trong (internal loop). Cấu
trúc bậc ba xác định cách những phần tử đó sắp xếp và tương tác
trong không gian. Phân cấp này giúp thu hẹp bài toán nhưng không làm nó đơn giản: hai
mô hình có thể dự đoán đúng phần lớn cặp base nitơ mà vẫn đặt sai tương quan giữa các miền
hoặc junction. RNA-Puzzles ghi nhận rằng một cấu trúc giữ đúng cấu trúc bậc hai vẫn có
thể lệch đáng kể khỏi cấu trúc thực nghiệm ở cấp ba chiều
\cite{rna_puzzles,rna_puzzles_toolkit}.

RNA còn có nhiều tương tác không chuẩn Watson--Crick, pseudoknot và chuyển động giữa
các trạng thái cấu dạng. Vì thế, trình tự giống nhau là bằng chứng có giá trị nhưng
không quyết định duy nhất cấu trúc; tương tự, một phân bố hình học phổ biến không thể thay
thế hoàn toàn thông tin về tiếp xúc xa và tổ chức miền. Các tổng quan gần đây thường
chia phương pháp dự đoán RNA 3D thành ba họ: mô phỏng dựa trên vật lý, lắp ráp dựa trên
tri thức/khuôn mẫu, và phương pháp học máy hoặc học sâu
\cite{rna_progress_review,state_of_rnart,computational_rna_review}.

\section{Biểu diễn cấu trúc và tính bất biến hình học}

Một nucleotide gồm base nitơ, đường ribose và nhóm phosphate. Nguyên tử \cprime{} nằm trong
đường ribose và cung cấp một điểm đại diện cho mỗi nucleotide. Khi chỉ giữ \cprime{},
một RNA dài $L$ trở thành ma trận $X\in\mathbb{R}^{L\times3}$. Đây là biểu diễn đầu ra
của cuộc thi và là mức biểu diễn chính của luận văn.

Tọa độ tuyệt đối không có ý nghĩa: cùng một cấu trúc có thể được xoay hoặc tịnh tiến mà
không thay đổi hình dạng. Do đó, phép so sánh cấu trúc phải bất biến đối với nhóm biến
đổi cứng trong không gian, thường ký hiệu là SE(3), gồm phép xoay và tịnh tiến. Kabsch
là thuật toán tìm phép xoay và tịnh tiến làm sai số bình phương nhỏ nhất giữa hai tập
điểm tương ứng. Phép chồng khít Kabsch, RMSD sau căn chỉnh và TM-score xử lý tính bất
biến này bằng cách tìm phép chồng khít cứng tối ưu. lDDT tránh căn chỉnh toàn cục bằng
cách so sánh trực tiếp các khoảng cách nội tại.

Biểu diễn \cprime{} làm giảm chi phí và phù hợp mục tiêu dự đoán nếp gấp, nhưng không
chứa hướng quay đầy đủ của base nitơ, liên kết hydro, xếp chồng base nitơ, cấu hình đường hoặc va
chạm toàn nguyên tử. Vì vậy một đánh giá hợp lý phải phân biệt ``độ chính xác đường
\cprime{}'' với ``tính hợp lý vật lý toàn nguyên tử''.

\section{Mô hình hóa dựa trên khuôn mẫu}

Mô hình hóa dựa trên khuôn mẫu (template-based modeling, TBM) tận dụng quan sát rằng các
RNA có quan hệ trình tự hoặc motif tương tự có thể chia sẻ một phần cấu trúc. Quy trình
TBM gồm bốn nhiệm vụ: nhận diện khuôn mẫu, căn chỉnh trình tự, chuyển thông tin hình học
và xây các vùng không được khuôn mẫu hỗ trợ. Các công cụ như ModeRNA và RNAComposer cho
thấy khả năng lắp ráp cấu trúc RNA từ khuôn mẫu hoặc đoạn cấu trúc đã biết
\cite{rna_composer,rna_puzzles}.

Điểm mạnh của TBM là truyền trực tiếp nếp gấp thực nghiệm, nhờ đó không phải tìm kiếm
toàn bộ không gian cấu dạng. Điểm yếu nằm ở chính giả định đó. Nếu thư viện không có
khuôn mẫu gần, nếu chỉ một motif ngắn tương đồng, hoặc nếu căn chỉnh sai tại insertion,
tọa độ được chuyển có thể thiếu hoặc sai. Các bộ đánh giá chuẩn RNA 3D cũng ghi nhận việc nhận
diện đúng khuôn mẫu là yếu tố quyết định đối với mô hình so sánh
\cite{state_of_rnart,rna_progress_review}.

\subsection{Truy hồi trình tự và nhận diện tương đồng xa}

Công cụ truy hồi nhanh thường dùng các đoạn con liên tiếp ($k$-mer) làm hạt giống để
loại phần lớn thư viện trước khi căn chỉnh chính xác. Lựa chọn $k$ tạo ra đánh đổi giữa
độ nhạy và chi phí: hạt giống dài giảm nhiễu nhưng dễ bỏ lỡ quan hệ xa. Vì cấu trúc có
thể được bảo tồn lâu hơn trình tự, chỉ một cơ chế tìm kiếm nhanh theo trình tự có thể
không đủ để phát hiện khuôn mẫu cấu trúc tiềm năng. Đây là lý do luận văn bổ sung một
điểm tương đồng tổng hợp, kết hợp căn chỉnh toàn cục, căn chỉnh cục bộ, thành phần trình
tự và mức trùng lặp của 3-mer.

\subsection{Căn chỉnh và chuyển tọa độ}

Căn chỉnh toàn cục phù hợp khi hai trình tự tương đồng trên phần lớn chiều dài; căn chỉnh
cục bộ tìm một miền hoặc motif bảo tồn. Kết quả căn chỉnh ánh xạ nucleotide thứ $i$ của
đối tượng sang nucleotide thứ $j$ của khuôn mẫu. Chỉ những vị trí có ánh xạ và có tọa
độ hợp lệ mới được sao chép trực tiếp. Những vị trí còn lại tạo thành khoảng trống và
cần được dựng bằng hình học, đoạn cấu trúc khác hoặc một mô hình dự đoán bổ sung. Sai lệch
ở bước căn chỉnh có thể truyền thành sai lệch hình học, vì vậy độ tương đồng trình tự,
độ phủ và độ đầy đủ tọa độ đều cần tham gia xếp hạng.

\section{Dự đoán dựa trên học sâu}

Sự phát triển của mô hình attention, biểu diễn ngôn ngữ và học hình học đã tạo ra nhiều
phương pháp RNA 3D mới. DRfold kết hợp học end-to-end về khung tọa độ cục bộ với các thế
năng hình học để hướng dẫn lắp ráp cấu trúc \cite{drfold}. RhoFold+ tận dụng một mô hình
ngôn ngữ RNA tiền huấn luyện trên tập trình tự lớn cùng thông tin tiến hóa để dự đoán
cấu trúc đơn chuỗi \cite{rhofoldplus}. DRfold2 thay đổi đáng kể kiến trúc tiền nhiệm bằng
mô hình ngôn ngữ RNA tổng hợp, module cấu trúc khử nhiễu và quy trình lựa chọn--tinh
chỉnh cấu dạng \cite{drfold2}. Boltz-1 là mô hình mở cho dự đoán cấu trúc và tương tác
phân tử, được dùng trong luận văn như một phương án khả thi cho RNA dài
\cite{boltz1}.

So với TBM, mô hình học sâu có thể sinh một giả thuyết ngay cả khi không có khuôn mẫu
gần trong thư viện truy vấn. Tuy nhiên, chúng vẫn học từ một tập cấu trúc thực nghiệm
hữu hạn. Độ chính xác có thể giảm khi chiều dài, họ RNA hoặc loại tương tác khác với
dữ liệu huấn luyện; hơn nữa một số phương pháp phụ thuộc vào MSA hoặc có giới hạn tài
nguyên. Vì vậy, kết quả từ mô hình học sâu cần được kiểm tra trùng lặp trình tự và mốc
thời gian huấn luyện trước khi diễn giải như bằng chứng tổng quát.

\section{Tinh chỉnh cấu trúc}

Cấu trúc sinh bởi TBM, lắp ráp đoạn cấu trúc hoặc học sâu có thể chứa khoảng cách bất
thường trên đường trục phân tử, va chạm, vùng thiếu hoặc lỗi hình học tại chỗ nối. Các phương pháp tinh chỉnh
dùng trường lực, thế năng thống kê, ràng buộc thực nghiệm hoặc tối ưu hình học để giảm
những lỗi đó. Các tổng quan về RNA 3D ghi nhận rằng refinement có thể sửa khuyết tật
cục bộ, nhưng mức cải thiện toàn cục phụ thuộc mạnh vào chất lượng cấu trúc ban đầu và
thông tin tiếp xúc xa \cite{rna_progress_review}.

Điểm cần phân biệt là ``làm cấu trúc hợp lý hơn theo hàm năng lượng'' và ``làm cấu trúc
gần thực nghiệm hơn''. Nếu một phương pháp được tối ưu trực tiếp để giảm va chạm hoặc
độ lệch góc, việc các đại lượng đó giảm chỉ chứng minh bộ tối ưu hoạt động theo mục tiêu
đã đặt. Để kết luận về độ chính xác, cần chỉ số độc lập như lDDT hoặc RMSD cục bộ. Đây
là nguyên tắc dẫn tới thiết kế đánh giá Geometry v2.

\section{Ước lượng chất lượng cấu trúc}

Sinh được một tập cấu trúc tốt chưa đồng nghĩa chọn được cấu trúc tốt. ARES cho thấy một
mạng học hình học có thể xếp hạng cấu trúc RNA từ tọa độ và cải thiện khả năng nhận diện
mô hình gần thực nghiệm \cite{ares}. Các hệ thống đánh giá RNA truyền thống và gần đây
cũng sử dụng thế năng thống kê, bản đồ tiếp xúc, đặc trưng nguyên tử hoặc dự đoán lDDT để
ước lượng chất lượng \cite{rnassess,ares}.

Trong pipeline cuối, độ tin cậy chỉ dùng để xếp hạng trong từng nguồn và điều chỉnh
cường độ Geometry; nó không được xem là xác suất đúng đã hiệu chỉnh giữa hai nguồn.
Một thí nghiệm khám phá về quality routing được trình bày ở phụ lục, nhưng kết quả đó
không tham gia contribution hay pipeline cuối.

\section{Độ tin cậy, hiệu chỉnh và cơ chế từ chối}

Độ tin cậy là đầu ra nội bộ của một mô hình, còn độ chính xác là quan hệ với cấu trúc
thực nghiệm. Hai khái niệm chỉ tương ứng chặt khi đã hiệu chỉnh. Giá trị 0.8 của TBM và
0.8 của DRfold2 không mặc nhiên cùng ý nghĩa; vì vậy pipeline không trộn hai thang điểm
thành một ranking toàn bank. Adaptive strength được ablate với fixed strength để kiểm
tra trực tiếp liệu việc dùng confidence có tạo giá trị hay không.

\section{Đánh giá chất lượng cấu trúc}

Không có một chỉ số duy nhất phản ánh đầy đủ mọi khía cạnh của RNA 3D. RNA-Puzzles và
các bộ công cụ đánh giá sử dụng nhiều chỉ số bổ sung cho nhau ở cấp độ toàn cục, cục bộ
và mạng tương tác \cite{rna_puzzles,rna_puzzles_toolkit,rnassess}. Vì luận văn chỉ có
đường \cprime{}, bốn nhóm chỉ số sau được chọn.

\subsection{TM-score và chất lượng nếp gấp tổng thể}

TM-score đo độ giống nếp gấp sau phép chồng khít cứng và chuẩn hóa theo chiều dài cấu
trúc tham chiếu \cite{tm_score}. Dạng trực giác:

\begin{equation}
\mathrm{TM} = \max_{R,t}\frac{1}{L}
\sum_{i=1}^{L}\frac{1}{1+(d_i/d_0(L))^2},
\end{equation}

trong đó $R,t$ là phép xoay và tịnh tiến, $d_i$ là khoảng cách sau căn chỉnh và $d_0(L)$
phụ thuộc chiều dài. Khoảng cách nhỏ đóng góp gần 1; khoảng cách lớn nhanh chóng bị giảm
trọng số. Luận văn dùng US-align để tính TM-score \cite{usalign}. Với năm cấu trúc có
TM lần lượt 0.42, 0.61, 0.38, 0.55 và 0.47, điểm của đối tượng là 0.61; ví dụ này giải
thích tại sao độ đa dạng của năm dự đoán có ý nghĩa.

\subsection{RMSD và sai lệch tọa độ}

Sau phép chồng khít cứng tối ưu:

\begin{equation}
\mathrm{RMSD}=\sqrt{\frac{1}{L}\sum_{i=1}^{L}
\left\|\hat{x}_i-y_i\right\|_2^2}.
\end{equation}

RMSD thấp hơn là tốt hơn, nhưng một miền đặt sai xa có thể chi phối toàn bộ giá trị.

\subsection{\texorpdfstring{\cprime{}-lDDT}{C1-prime-lDDT} và độ chính xác cục bộ}

lDDT đo khả năng bảo toàn các khoảng cách lân cận mà không cần căn chỉnh toàn cục
\cite{lddt}. Bản thích nghi của luận văn dùng một atom \cprime{} mỗi nucleotide, bán
kính xét cặp 15~\angstrom{} và các ngưỡng sai lệch 0.5, 1, 2 và 4~\angstrom{}. Do không
dùng đầy đủ nguyên tử, tên chính xác là \cprime{}-lDDT, không phải lDDT toàn nguyên tử.

\subsection{RMSD theo cửa sổ trượt}

RMSD được tính lại sau khi căn chỉnh riêng từng đoạn liên tiếp dài 9, 15 và 31
nucleotide. Cửa sổ ngắn nhạy với hình học tại chỗ; cửa sổ dài hơn bắt đầu phản ánh tổ
chức rộng. Ba kích thước này là quy mô phân tích được xác định trước, không phải hằng số
sinh học phổ quát.

\section{Khoảng trống mà luận văn giải quyết}

Các công trình trước cho thấy TBM mạnh khi nhận đúng khuôn, còn học sâu có thể tạo nếp
gấp bổ sung. Khoảng trống của luận văn là thiếu phân rã white-box dưới kiểm soát thời
gian và thiếu bằng chứng độc lập rằng một refiner hình học coarse-grained làm cấu trúc
gần native hơn, thay vì chỉ giảm chính loss đã đặt. Geometry v2 vì vậy được đánh giá
bằng lDDT, RMSD cửa sổ và TM trên split 60/20/20, kèm đối chứng và component ablation.

% -----------------------------------------------------------------------------
% CHAPTER 3
% -----------------------------------------------------------------------------
\chapter{Dữ liệu và quy trình đánh giá}

\section{Các nhóm dữ liệu}

\begin{longtable}{p{0.24\textwidth}p{0.29\textwidth}p{0.36\textwidth}}
\toprule
Nhóm & Nội dung & Vai trò\\
\midrule
\endhead
Bảng trình tự & mã đối tượng, trình tự, mô tả, mốc thời gian & Đầu vào và siêu dữ liệu cho các tập huấn luyện, hiệu chỉnh và đánh giá.\\
Bảng nhãn & tọa độ thực nghiệm \cprime{} & Học prior trên 60 RNA; calibration tham số; đánh giá một lần sau khi đóng băng phương pháp.\\
\texttt{PDB\_RNA/*.cif} & chuỗi RNA, nguyên tử, nucleotide biến đổi, ngày công bố & Thư viện khuôn mẫu chính.\\
MSA & các trình tự họ hàng đã căn chỉnh & Đầu vào tiềm năng cho mô hình học sâu; nhánh TBM hiện tại không phụ thuộc MSA.\\
Dữ liệu mô hình học sâu & bộ trọng số đã huấn luyện và cấu trúc DRfold2/Boltz & Nguồn giả thuyết nếp gấp bổ sung.\\
\bottomrule
\end{longtable}

\section{Thư viện khuôn mẫu}

Phiên bản dữ liệu được phân tích bằng Gemmi gồm 8,670 tệp CIF, tạo 23,869 chuỗi RNA
hoặc chuỗi lai RNA--DNA với khoảng 10.86 triệu nucleotide. Khoảng 99.9\% nucleotide có
\cprime{} được mô hình hóa. Nucleotide biến đổi được ánh xạ về A/C/G/U khi có quy tắc;
tọa độ thiếu được giữ là NaN để bước dựng khoảng trống xử lý, thay vì nội suy âm thầm
trong lúc đọc dữ liệu.

Hai thư viện tìm kiếm được xây:

\begin{itemize}
  \item khoảng 20,844 chuỗi hợp lệ trong cơ sở dữ liệu MMseqs2;
  \item 7,155 trình tự duy nhất trong thư viện tìm kiếm mở rộng, giữ đại diện tọa độ
  có ngày công bố sớm nhất.
\end{itemize}

\section{Tập phát triển CASP15}

Tập phát triển cục bộ gồm 12 đối tượng: R1107, R1108, R1116, R1117v2, R1126, R1128,
R1136, R1138, R1149, R1156, R1189 và R1190, chiều dài từ 30 đến 720 nucleotide. Tập
này được dùng để đánh giá việc tìm khuôn mẫu, độ đa dạng của nguồn dự đoán và Geometry
v2. Vì chỉ có 12 RNA độc lập, thống kê được lấy mẫu lại theo RNA, không coi hàng nghìn
nucleotide trong cùng phân tử là các mẫu độc lập.

\section{Tập 60/20/20 cho Geometry v2 và phân tích lai}

Một trăm RNA dài tối đa 100 nucleotide được sắp theo thời gian và gom nhóm họ trình tự.
Tập này không còn phục vụ mô hình GeoFuse trong mạch chính; nó được chuyển đúng vai trò
train--calibration--validation cho Geometry v2:

\begin{table}[H]
\centering
\caption{Phân chia thời gian/họ cho đánh giá Geometry v2}
\begin{tabular}{lrrlll}
\toprule
Tập & Số RNA & Số họ & Khoảng ngày & TBM & DRfold2\\
\midrule
Huấn luyện prior & 60 & 54 & 2024-01-10--2024-12-04 & 180 & 120\\
Hiệu chỉnh & 20 & 17 & 2024-12-11--2025-01-08 & 60 & 40\\
Xác nhận mới nhất & 20 & 16 & 2025-01-15--2025-03-26 & 60 & 40\\
\bottomrule
\end{tabular}
\end{table}

Không có họ hoặc nhóm trình tự trùng hoàn toàn xuất hiện ở nhiều tập. Mỗi RNA có đúng ba
ứng viên TBM và hai DRfold2 hữu hạn. Sáu mươi RNA đầu chỉ ước lượng prior hình học; 20
RNA hiệu chỉnh chọn cấu hình trong một danh sách khai báo trước; cấu hình và checksum
được khóa trước khi mở 20 RNA xác nhận. Một đối tượng lỗi kỹ thuật được thay theo thứ tự
dự phòng định trước, không dựa trên độ chính xác.

\section{Kiểm tra trùng lặp và mốc huấn luyện của mô hình học sâu}

R1128 có trình tự trùng hoàn toàn với một mục trong FASTA huấn luyện DRfold2 và tạo
TM-score rất cao. Vì vậy, kết quả về khả năng bổ sung nếp gấp được báo cáo theo hai cách:
đủ 12 đối tượng theo thiết lập cuộc thi và phân tích độ nhạy loại R1128. R1138 dùng Boltz
có mốc dữ liệu huấn luyện được công bố sớm hơn ngày phát hành cấu trúc đích. Việc kiểm
tra này giảm nguy cơ diễn giải ghi nhớ dữ liệu thành khả năng khái quát.

\section{Kiểm soát rò rỉ theo thời gian}

Với R1107 có mốc 2022-05-28, cấu trúc 7QR4 phát hành ngày 2022-10-26 phải bị loại khỏi
thư viện khuôn mẫu dù có trình tự rất gần. Một phép tái tạo cho phép rò rỉ đạt TM 0.9355,
trong khi phiên bản an toàn theo thời gian đạt 0.2983. Chênh lệch này cho thấy bộ lọc
ngày công bố là điều kiện trung tâm để kết quả có giá trị khoa học.

Ngoài ngày công bố, quy trình còn loại PDB của chính đối tượng, kiểm tra trùng trình tự,
gom nhóm theo họ và xem xét mốc dữ liệu huấn luyện của mô hình tiền huấn luyện. Những
kiểm tra này không bảo đảm tuyệt đối rằng mọi quan hệ tiến hóa xa đã được loại, nhưng
giảm đáng kể các hình thức rò rỉ có thể quan sát và tái lập.

\section{Quy trình sử dụng nhãn}

Cấu trúc thực nghiệm xuất hiện ở ba vai trò tách biệt. Nhãn của 60 RNA đầu chỉ dùng để
ước lượng phân bố khoảng cách, kích thước, góc và góc xoắn. Tập hiệu chỉnh được dùng để
chọn source restraint, adaptive strength, learning rate và số bước trong các giá trị
khai báo trước. Tập xác nhận mới nhất chỉ được mở sau khi ghi cấu hình đóng băng; tham
số không được thay đổi theo kết quả này. Đối với Kaggle, cấu trúc thực nghiệm của tập
ẩn hoàn toàn không có sẵn.

Khái niệm ngoài nếp gấp (out-of-fold, OOF) bảo đảm ứng viên của một RNA không được tạo
bởi một mô hình huấn luyện trên chính RNA hoặc họ của nó. Ngoài OOF theo mô hình, thứ tự
thời gian giữa ba tập tạo một phép kiểm tra triển vọng: prior luôn đến trước calibration,
và calibration luôn đến trước validation.

\section{Chỉ số và đơn vị phân tích}

TM-score tốt nhất trong năm cấu trúc là chỉ số chính cho chất lượng nếp gấp. Geometry
v2 được đánh giá thêm bằng \cprime{} RMSD, \cprime{}-lDDT và RMSD cửa sổ 9, 15, 31
nucleotide. Clash proxy, sharp-kink, bán kính quán tính và NLL góc/góc xoắn chỉ là
diagnostic của objective; chúng không thay thế các chỉ số native độc lập.

Mỗi RNA có trọng số bằng nhau trong giá trị trung bình. So sánh theo cặp sử dụng
lấy mẫu lại 10,000 lần ở cấp RNA với hạt giống ngẫu nhiên 2025 để tạo khoảng tin cậy
95\%. Cách làm này
tránh ngụy lặp, tức là coi nhiều nucleotide tương quan trong một RNA như nhiều thí
nghiệm độc lập.

% -----------------------------------------------------------------------------
% CHAPTER 4
% -----------------------------------------------------------------------------
\chapter{Khung dự đoán lai đề xuất}

\section{Tổng quan}

\begin{figure}[H]
\centering
\fbox{\begin{minipage}{0.91\textwidth}
\centering
\textbf{Đầu vào: trình tự + mốc thời gian}\\[0.2cm]
$\downarrow$\\
\begin{tabular}{c@{\hspace{1.3cm}}c}
\textbf{Dựa trên khuôn mẫu} & \textbf{Dựa trên học sâu}\\
MMseqs2 + tìm kiếm mở rộng & DRfold2 / Boltz cho RNA dài\\
lọc $\rightarrow$ căn chỉnh $\rightarrow$ chuyển tọa độ & bộ trọng số cố định $\rightarrow$ cấu trúc\\
\end{tabular}\\[0.25cm]
$\downarrow$\\
Tập cấu trúc ban đầu: T1, T2, T3, P1, P2\\
$\downarrow$\\
Tinh chỉnh bảo thủ Geometry v2 trên từng ứng viên\\
$\downarrow$\\
Phân bổ cố định 3 TBM + 2 DRfold2 (fallback về TBM nếu cần)\\
$\downarrow$\\
\textbf{Năm cấu trúc \cprime{} và tệp \texttt{submission.csv}}
\end{minipage}}
\caption{Khung dự đoán tổng thể. Cấu trúc thực nghiệm chỉ xuất hiện trong huấn luyện và đánh giá, không xuất hiện khi suy luận.}
\label{fig:pipeline}
\end{figure}

\section{Sinh cấu trúc từ khuôn mẫu}

\subsection{Truy hồi nhanh bằng MMseqs2}

Trình tự cần dự đoán được truy vấn trong khoảng 20,844 chuỗi bằng MMseqs2 với $k=13$.
Nhánh này nhanh và hiệu quả khi có quan hệ tương đồng rõ. Kết quả ở giai đoạn này chỉ
là danh sách khuôn mẫu tiềm năng và căn chỉnh sơ bộ; chưa có tọa độ nào được sao chép.

\subsection{Tìm kiếm tương đồng mở rộng}

Song song với MMseqs2, phương pháp quét 7,155 trình tự khuôn mẫu duy nhất. Điểm tương
đồng tổng hợp được tính bởi:

\begin{equation}
C(S,T)=0.4G(S,T)+0.3L(S,T)+0.2F(S,T)+0.1K_3(S,T),
\end{equation}

với $G$ là độ tương đồng căn chỉnh toàn cục, $L$ là độ tương đồng căn chỉnh cục bộ,
$F$ tổng hợp thành phần base nitơ, hàm lượng GC, cặp nucleotide liên tiếp, entropy, đoạn
lặp và chiều dài;
$K_3$ là tỷ lệ 3-mer chung. ``Mở rộng'' ở đây nghĩa là đánh giá toàn bộ thư viện trình
tự hợp lệ thay vì chỉ các kết quả MMseqs2; không phải vét cạn không gian cấu hình 3D.

\subsection{Hợp nhất kết quả, lọc và xếp hạng}

Hai danh sách kết quả được hợp nhất. Mỗi khuôn mẫu phải thỏa:

\begin{equation}
t_{release}(T)<t_{cutoff}(S),
\end{equation}

và không thuộc PDB tham chiếu của chính đối tượng. Các cặp được căn chỉnh lại toàn cục
và xếp hạng theo tích của mức đồng nhất trình tự, độ phủ đối tượng và độ đầy đủ tọa độ
khuôn mẫu. Phương pháp ưu tiên các PDB khác nhau để tránh năm dự đoán cùng một nguồn.

\subsection{Chuyển tọa độ}

Phép căn chỉnh tạo ánh xạ $m(i)$ từ nucleotide $i$ của đối tượng tới nucleotide $j$
của khuôn mẫu. Nếu $m(i)=j$ và khuôn mẫu có \cprime{} hữu hạn:

\begin{equation}
x_i^{TBM}=x_j^{template},\qquad support_i=1.
\end{equation}

Nếu nucleotide không được ánh xạ, $support_i=0$ và tọa độ được dựng ở bước lấp khoảng
trống. Mặt nạ hỗ trợ phân biệt tọa độ sao chép trực tiếp với tọa độ suy dựng.

\subsection{Lấp khoảng trống và gán độ tin cậy}

Khoảng trống ngắn ở giữa được nội suy giữa hai điểm neo. Khoảng trống dài kết hợp nội
suy với độ cong vuông góc để tránh đường thẳng co cụm. Phần thiếu ở đầu hoặc cuối chuỗi
được kéo dài theo hướng backbone gần nhất. Nếu không có khuôn mẫu, một phương án dựng
thô sinh nhiều đường \cprime{} dựa trên thân xoắn và vòng. Nucleotide được chuyển trực tiếp
có độ tin cậy cao hơn; độ tin cậy trong vùng suy dựng giảm theo khoảng cách tới điểm neo.

\section{Sinh cấu trúc bằng mô hình học sâu}

\subsection{DRfold2}

Chương trình ghi trình tự thành FASTA, chạy 20 bộ trọng số đã huấn luyện của cấu hình \texttt{cfg97},
dự đoán các ràng buộc/thế năng hình học, dùng PotentialFold tối ưu cấu trúc thô và Arena
tạo PDB trước khi trích \cprime{}. Nhiều bộ trọng số tạo nhiều giả thuyết. Các cấu trúc
được xếp hạng theo tiêu chí nội bộ của mô hình, nhưng độ tin cậy thô không được xem là
độ chính xác so với thực nghiệm.

\subsection{Boltz cho RNA dài}

Boltz được kiểm tra cho R1138 dài 720 nucleotide khi DRfold2 vượt giới hạn GPU 16 GB.
Do tập xác nhận giới hạn chiều dài $\le100$, DRfold2 được dùng thống nhất trong thí
nghiệm lựa chọn nguồn. Boltz vì vậy cung cấp bằng chứng khả thi và phương án cho RNA
dài, không phải bộ sinh mặc định trong thí nghiệm xác nhận cuối.

\subsection{Biểu diễn ứng viên thống nhất}

Hai nguồn được chuẩn hóa thành một cấu trúc dữ liệu chung:

\begin{lstlisting}
target_id, sequence
coords:        [L, 3]
confidence:    [L]
support_mask:  [L]
source, model, checkpoint, template_id, cutoff provenance
\end{lstlisting}

Biểu diễn này cho phép Geometry v2 xử lý hai nguồn bằng cùng một giao diện, đồng thời
giữ nguồn gốc và phân biệt nucleotide được khuôn hỗ trợ với nucleotide do gap-fill dựng.

\section{Geometry v2}

\subsection{Mục tiêu thiết kế}

Geometry v2 không dự đoán lại nếp gấp từ đầu. Phương pháp tìm cấu trúc $X$ có hình học
\cprime{} cục bộ hợp lý hơn nhưng vẫn gần cấu trúc nguồn $X_0$. Cấu trúc ban đầu luôn
được giữ trong tập ứng viên để tinh chỉnh không loại bỏ một dự đoán vốn đã tốt.

\subsection{Các phân bố hình học thực nghiệm}

Từ các cấu trúc thực nghiệm cũ hợp lệ theo thời gian, phương pháp đo:

\begin{itemize}
  \item khoảng cách \cprime{} liên tiếp;
  \item pseudo-angle từ ba \cprime{} liên tiếp;
  \item góc xoắn giả có dấu từ bốn \cprime{} liên tiếp;
  \item bán kính quán tính theo chiều dài;
  \item khoảng cách tới láng giềng không kề nhau, chỉ dùng để ước lượng ngưỡng clash;
  \item chỉ dấu pair-like, chỉ dùng để chọn prior angle/torsion theo ngữ cảnh.
\end{itemize}

Góc phẳng và góc xoắn được chia thành 72 khoảng, làm trơn bằng nhân có trọng số
\begin{equation}
(1,2,3,4,3,2,1),
\end{equation}
thêm
một lượng đếm nhỏ để tránh xác suất bằng 0 rồi đổi sang âm log-xác suất (NLL). NLL thấp
nghĩa là hình học thuộc vùng thường gặp hơn trong dữ liệu dùng để xây phân bố.

Chỉ dấu ``có dạng cặp'' (pair-like) là một đại lượng gần đúng chỉ dựa trên thông tin có
ở thời điểm suy luận. Một cặp được cân nhắc nếu base nitơ thuộc AU, UA, GC, CG, GU hoặc
UG, cách nhau ít nhất bốn vị
trí trên trình tự và có khoảng cách \cprime{} quanh
$10.5\pm2.5$~\angstrom. Các cặp được chọn tuần tự theo độ gần mục tiêu, mỗi nucleotide
tối đa một cặp; các cặp giao nhau được phép để không loại nút giả (pseudoknot). Chỉ dấu
này không phải nhãn cặp base nitơ thực nghiệm và không tạo một loss pair-like riêng;
nó chỉ chọn histogram angle/torsion tương ứng cho từng vị trí.

\subsection{Hàm mục tiêu}

\begin{align}
\mathcal{L}(X)={}&1.5\mathcal{L}_{source}
+s\big(1.0\mathcal{L}_{backbone}
+0.3\mathcal{L}_{clash}
+0.02\mathcal{L}_{Rg}\nonumber\\
&\qquad +0.30\mathcal{L}_{angle}
+0.15\mathcal{L}_{torsion}
+20.0\mathcal{L}_{kink}\big),\\
s={}&0.2+0.8(1-c_{global}).
\end{align}

$\mathcal{L}_{source}$ là khoảng cách Huber về $X_0$, có trọng số theo độ tin cậy.
$\mathcal{L}_{backbone}$ phạt khoảng cách nucleotide liên tiếp lệch khỏi phân bố thực nghiệm.
$\mathcal{L}_{clash}$ phạt \cprime{} không kề nhau quá gần.
$\mathcal{L}_{Rg}$ giữ kích thước tổng thể hợp lý theo chiều dài.
$\mathcal{L}_{angle}$ và $\mathcal{L}_{torsion}$ phạt vùng có âm log-xác suất (NLL) cao.
$\mathcal{L}_{kink}$ ngăn tạo hoặc làm xấu góc gập sắc.

Trọng số source 1.5 được chọn trên 20 RNA hiệu chỉnh từ các giá trị 1.5, 3.0 và 6.0;
mọi thành phần còn lại giữ cấu hình khai báo trước. Một góc gập sắc được ghi nhận khi pseudo-angle nhỏ hơn $70^\circ$. Bộ tối ưu dùng biên
bảo vệ $75^\circ$: nếu góc của cấu trúc nguồn vốn lớn hơn biên, góc sau tinh chỉnh không
nên rơi xuống dưới; nếu góc nguồn đã nhỏ, cận dưới là chính góc đó để không làm gập sắc
hơn. Tọa độ được tối ưu bằng Adam trong 300 bước, tốc độ học 0.04 và giới hạn chuẩn
gradient ở 10.

\iffalse
\section{GeoFuse}

\subsection{Bài toán lựa chọn nguồn}

Ý tưởng giữ vùng tốt từ TBM và thay vùng còn lại bằng mô hình học sâu chỉ có ý nghĩa
nếu có thể ước lượng vùng nào tốt mà không nhìn cấu trúc thực nghiệm. GeoFuse tách vấn
đề thành hai tầng: trước hết học một mô hình phân loại nguồn tốt hơn theo nucleotide;
sau đó áp dụng một quy tắc bảo thủ để quyết định có tạo cấu trúc hợp nhất hay không.

\subsection{Phân cụm theo nếp gấp tổng thể}

Phương pháp tính self-TM đối xứng giữa mọi cặp cấu trúc, chuyển thành khoảng cách
$d=1-\mathrm{TM}$ và phân cụm phân cấp theo liên kết đầy đủ (complete-link): khoảng
cách giữa hai cụm được lấy bằng cặp phần tử xa nhau nhất. Ngưỡng 0.35 được chọn trên
tập hiệu chỉnh. Chỉ một cấu trúc khuôn mẫu và một cấu trúc học sâu trong cùng cụm mới
được ghép cặp. Điều này tránh lấy trung bình hai nếp gấp tổng thể khác nhau.

\subsection{Chồng khít bền vững trước ngoại lệ}

Cấu trúc học sâu được chồng khít cứng lên cấu trúc khuôn mẫu bằng phép xoay và tịnh tiến
Kabsch. Vị trí dùng để khớp phải hữu hạn ở hai bên, được khuôn mẫu hỗ trợ và có độ tin
cậy khuôn mẫu ít nhất 0.25. Sau mỗi vòng, 80\% sai số dư thấp nhất được giữ lại và 20\%
điểm ngoại lệ tạm thời bị loại; quá trình lặp ba lần. Phép biến đổi không kéo giãn hay
bẻ cong cấu trúc.

Sau khi chồng khít:

\begin{equation}
d_i=\left\|\widehat{x}_i^{P}-x_i^{T}\right\|_2,
\end{equation}

Độ bất đồng $d_i$ được làm trơn bằng cửa sổ 15 nucleotide. Đại lượng này chỉ cho biết
hai nguồn khác nhau bao nhiêu, không cho biết nguồn nào đúng.

\subsection{Đặc trưng và mô hình lựa chọn nguồn}

Mỗi nucleotide có 22 đặc trưng: độ tin cậy của hai nguồn; thứ hạng độ tin cậy trong
chuỗi; trạng thái hỗ trợ của khuôn mẫu và độ sâu trong vùng được hỗ trợ; độ bất đồng giữa
hai nguồn; năm đặc trưng hình học cục bộ cho mỗi nguồn (độ lệch đường trục, NLL góc, NLL
góc xoắn, góc gập sắc và chỉ dấu có dạng cặp); mã hóa nhị phân một-trong-bốn cho A/C/G/U;
và logarit chiều dài chuỗi.

Đặc trưng được chuẩn hóa bằng trung bình và độ lệch chuẩn của tập huấn luyện. Mô hình
gồm hai lớp tích chập một chiều (Conv1D) 32 kênh, nhân rộng 5 vị trí, hàm kích hoạt
SiLU và tỷ lệ bỏ nút (dropout) 0.05, sau đó một lớp Conv1D trả một điểm chưa chuẩn hóa
(logit). Hai lớp với nhân rộng 5 vị trí cho vùng quan sát xấp xỉ 9 nucleotide. Hàm
sigmoid biến điểm này về khoảng từ 0 đến 1 và trả:

\begin{equation}
p_i=P(\text{DRfold2 tốt hơn TBM tại nucleotide }i).
\end{equation}

Trong huấn luyện, nhãn chính so sánh \cprime{}-lDDT theo nucleotide của hai nguồn với
cấu trúc thực nghiệm. Khi suy luận, cấu trúc thực nghiệm và mọi đặc trưng suy ra từ nhãn
hoàn toàn vắng mặt.

\subsection{Hợp nhất có chọn lọc và cơ chế từ chối}

Tập hiệu chỉnh chọn ngưỡng quyết định 0.75; biên an toàn 0.15 tạo ngưỡng thực tế
$p_i\ge0.90$. Một nucleotide chỉ được chọn nếu đồng thời:

\begin{equation}
p_i\ge0.90,\qquad 0.5\le d_i\le12~\angstrom.
\end{equation}

Phải tồn tại ít nhất 7 nucleotide liên tục. Mặt nạ nhị phân được làm trơn bằng bộ lọc
Gauss, tạo
$\alpha_i\in[0,1]$:

\begin{equation}
x_i^{fused}=(1-\alpha_i)x_i^{T}+\alpha_i\widehat{x}_i^{P}.
\end{equation}

Nếu không có đoạn liên tục đạt điều kiện, phương pháp không tạo cấu trúc hợp nhất. Đây
là cơ chế từ chối hợp nhất có chủ đích. Các cấu trúc cha luôn còn trong tập ứng viên.
Geometry v2 có thể tinh chỉnh cấu trúc hợp nhất để giảm sai lệch tại đường nối nhưng
không thay thế cấu trúc cha.

\section{Lựa chọn năm cấu trúc cuối}

Tập ứng viên gồm cấu trúc TBM và học sâu ban đầu, các cấu trúc sau Geometry v2 và cấu
trúc hợp nhất nếu có. Bộ chọn sử dụng điểm chất lượng không dựa trên cấu trúc thực
nghiệm, với trọng số độ tin cậy nguồn 0.25, tỷ lệ được khuôn mẫu hỗ trợ 0.20, tỷ lệ
chỉ dấu có dạng cặp 0.10, NLL góc phẳng 0.15, NLL góc xoắn 0.10, số va chạm trên
nucleotide 0.07, độ lệch đường trục 0.08 và góc gập sắc 0.05. Chỉ số có chiều
``thấp hơn là tốt hơn'' được
đảo thứ hạng trước khi cộng.

Quá trình lựa chọn còn có trọng số đa dạng 0.25, thưởng cho cụm nếp gấp mới 0.20 và mức
ủng hộ của cụm 0.10. Mục tiêu là tránh năm cấu trúc đều thuộc cùng một nếp gấp hoặc một
nguồn. TM-score thực nghiệm chỉ được tính sau khi danh sách năm cấu trúc đã được cố định.

\section{Ghi tệp nộp bài}

Mỗi nucleotide là một dòng:

\begin{lstlisting}
ID,resname,resid,x_1,y_1,z_1,...,x_5,y_5,z_5
\end{lstlisting}

Các bước kiểm tra gồm đúng ID và thứ tự mẫu, nucleotide khớp trình tự, đủ năm cấu dạng,
kích thước mảng đúng, tọa độ hữu hạn và không có ID trùng.

\section{Mã giả đầy đủ}

\begin{lstlisting}
for target in test_sequences:
    hits = mmseqs_search(target.sequence, k=13)
    hits += exhaustive_composite_search(target.sequence)
    hits = temporal_and_self_filter(hits, target)

    tbm = []
    for template in realign_rank_distinct(hits):
        coords, support = transfer_C1prime(template, target)
        coords, confidence = geometry_gap_fill(coords, support)
        tbm.append(candidate(coords, confidence, support))
    tbm = top3_distinct_PDB(tbm)

    pretrained = DRfold2_frozen(target.sequence, top2=True)
    raw = normalize(tbm + pretrained)
    augmented = raw + geometry_v2(raw)

    clusters = self_TM_complete_link(raw, threshold=0.35)
    for template, pretrained in same_fold_mixed_pairs(clusters):
        P = quality_router(template, pretrained)
        segment = contiguous_run(
            (P >= 0.90) & (disagreement >= 0.5)
                        & (disagreement <= 12.0),
            minimum_length=7)
        if segment exists:
            fused = gaussian_smooth_fusion(template, pretrained, segment)
            augmented += [fused, geometry_v2(fused)]
        else:
            abstain

    final5 = native_blind_quality_diversity_select(augmented)
    write_submission_rows(target, final5)
\end{lstlisting}

\fi

\section{Phân bổ năm cấu trúc cuối và ghi tệp}

Pipeline cuối không ghép tọa độ giữa hai nguồn. Với RNA không quá 600 nucleotide, hệ
thống giữ ba ứng viên TBM xếp hạng cao nhất và hai ứng viên DRfold2; nếu DRfold2 lỗi
hoặc vượt giới hạn tài nguyên, vị trí đó quay về TBM. Geometry v2 tinh chỉnh từng ứng
viên độc lập. Phân bổ cố định này không cần nhãn native khi suy luận và phản ánh đúng
notebook Kaggle đã được chấm.

Mỗi nucleotide là một dòng \texttt{ID,resname,resid,x\_1,y\_1,z\_1,...,x\_5,y\_5,z\_5}.
Kiểm tra cuối bảo đảm đúng ID/thứ tự mẫu, đúng nucleotide, đủ năm cấu dạng, tọa độ hữu
hạn và không có ID trùng.

\begin{lstlisting}
for target in test_sequences:
    hits = mmseqs_search(target.sequence, k=13)
    hits += exhaustive_composite_search(target.sequence)
    hits = temporal_and_self_filter(hits, target)
    tbm = transfer_rank_distinct_and_gap_fill(hits, top=3)
    deep = DRfold2_frozen(target.sequence, top=2)
    raw5 = tbm + deep if deep succeeds else tbm_fallback_to_five
    final5 = [geometry_v2(candidate) for candidate in raw5]
    validate_and_write_submission(target, final5)
\end{lstlisting}

% -----------------------------------------------------------------------------
% CHAPTER 5
% -----------------------------------------------------------------------------
\chapter{Thiết kế thí nghiệm}

\section{Chiến lược kiểm định giả thuyết}

Các thí nghiệm đi theo chuỗi component--hypothesis--ablation--conclusion. TBM được tách
thành retrieval source, bốn thành phần composite score, reranking, distinct-PDB và gap
filling. Tính bổ sung TBM--DRfold2 được kiểm tra bằng allocation giữ cố định số candidate.
Geometry v2 được mổ bằng cumulative/leave-one-out, đối chứng đơn giản, sensitivity,
convergence và cuối cùng là xác nhận 20 RNA mới nhất.

Prior được ước lượng trên 60 RNA; tham số Geometry chỉ được chọn trên 20 hiệu chỉnh từ
danh sách khai báo trước. Tập xác nhận mới nhất được mở sau khi ghi tệp cấu hình đóng
băng. Native reference cho mỗi candidate refined được cố định theo candidate thô để
phương pháp không được lợi do đổi sang cấu dạng native dễ hơn.

\section{Liên hệ giữa câu hỏi nghiên cứu và thí nghiệm}

\begin{longtable}{p{0.08\textwidth}p{0.29\textwidth}p{0.27\textwidth}p{0.24\textwidth}}
\toprule
ID & Giả thuyết & So sánh & Chỉ số chính\\
\midrule
\endhead
E0 & TBM phải an toàn theo thời gian & safe so với cho phép template tương lai & TM, availability\\
E1 & Composite tăng recall & MMseqs2, composite và hợp hai nguồn & availability, top-1/best-5 TM, coverage\\
E2 & Reranking/gap-fill có vai trò riêng & component score; identity; identity$\times$coverage; distinct-PDB; linear/current gap-fill & top-1/3/5 TM, SW-RMSD9\\
E3 & DRfold2 bổ sung TBM & 3T; 2T+1D; 1T+2D; 3T+2D, Geometry off/on & best-of-$N$ TM, oracle, bootstrap\\
E4 & Thành phần Geometry có tác dụng gì & cumulative, leave-one-out, fixed/adaptive, source/lr/step sensitivity & TM, lDDT, SW-RMSD, diagnostic\\
E5 & Geometry vượt đối chứng đơn giản & raw, source--backbone, full Geometry & chỉ số native độc lập\\
E6 & Geometry khái quát khi nào & 20 RNA newest; stratify source/quality/length/support & paired target bootstrap\\
\bottomrule
\end{longtable}

\section{Phương pháp đối chứng}

Đối chứng TBM gồm MMseqs2 đơn lẻ, từng subset của composite score, identity-only ranking
và gap-fill tuyến tính. Đối chứng Geometry ``ngu'' nhưng công bằng chỉ có source restraint
và backbone loss, cùng số bước và candidate. Full Geometry được so với fixed strength,
leave-one-out từng loss và raw không tinh chỉnh. Diagnostic được báo cáo để giải thích
cơ chế nhưng gate chỉ dùng TM, \cprime{}-lDDT và sliding-window RMSD độc lập.

\section{Điểm lý tưởng và điểm lựa chọn}

Điểm lý tưởng (oracle) chọn cấu trúc tốt nhất bằng cách nhìn cấu trúc thực nghiệm sau
khi đánh giá; nó chỉ đo trần chất lượng của tập ứng viên. Điểm lựa chọn được tạo bằng quy tắc không sử
dụng cấu trúc thực nghiệm và có thể áp dụng cho dữ liệu mới. Chênh lệch giữa điểm lý
tưởng và điểm lựa chọn là phần chất lượng còn mất ở khâu xếp hạng. Khi điểm lý tưởng
tăng nhưng điểm lựa chọn không tăng, hệ thống đã sinh được cấu trúc tốt hơn nhưng chưa
nhận ra chúng một cách đáng tin.

\section{Phân tích thống kê}

Mọi RNA có trọng số bằng nhau bất kể chiều dài. So sánh theo cặp dùng phép lấy mẫu lại
10,000 lần ở cấp RNA với hạt giống ngẫu nhiên 2025 để tạo khoảng tin cậy 95\%.
Nucleotide trong cùng một RNA
không được coi là đối tượng độc lập. Khi khoảng tin cậy của chênh lệch chứa 0, kết quả
được diễn giải là chưa đủ bằng chứng về một thay đổi ổn định.

\section{Tiêu chí đánh giá Geometry v2}

Trước khi mở kết quả xác nhận, tiêu chí được đặt là \cprime{}-lDDT phải tăng, RMSD phải
cải thiện ở ít nhất hai trong ba quy mô cửa sổ và chênh lệch TM không thấp hơn $-0.005$.
Tiêu chí này buộc phương pháp thể hiện lợi ích bằng chỉ số độc lập đồng thời hạn chế làm
hỏng nếp gấp tổng thể.

\section{Factorial và giới hạn allocation}

Cache confirmatory có ba TBM và hai DRfold2 trên mỗi RNA. Factorial chính vì vậy giữ
cố định $N=3$: 3T+0D so với 2T+1D, mỗi bank có Geometry off/on. Bảng allocation còn
báo cáo mọi tổ hợp khả thi ở $N=1,2,3$ và production bank 3T+2D. Không có dữ liệu để
giả lập trung thực 5T+0D hay 0T+5D; luận văn ghi đây là giới hạn thay vì đệm bản sao.

\section{Môi trường tính toán và tái lập}

Hệ thống được phát triển trong WSL với môi trường conda \texttt{rna-fold}; mã nguồn và
cấu hình thí nghiệm được quản lý bằng Git. Suy luận DRfold2 sử dụng GPU Kaggle khi tài
nguyên cục bộ không đủ. Các tệp cấu trúc và bộ trọng số lớn được lưu ngoài Git, trong
khi tệp kê khai, nguồn gốc dữ liệu, cấu hình, báo cáo và bảng tổng hợp được lưu trong
kho mã. \TODO{Bổ sung CPU, GPU, RAM, phiên bản WSL, CUDA, Python, PyTorch và thời gian
chạy cuối từ tệp kê khai.}

% -----------------------------------------------------------------------------
% CHAPTER 6
% -----------------------------------------------------------------------------
\chapter{Kết quả}

\section{RQ1: Mở rộng khả năng tìm khuôn mẫu}

Giả thuyết đầu tiên cho rằng một cơ chế truy hồi nhanh có thể bỏ lỡ những khuôn mẫu xa,
và việc bổ sung tìm kiếm tương đồng mở rộng sẽ cải thiện chất lượng tập dự đoán. Hai
phương pháp được so sánh trên cùng 12 đối tượng CASP15, cùng bộ lọc thời gian, cùng bước
chuyển tọa độ và cùng bộ tinh chỉnh; khác biệt duy nhất là việc bật tìm kiếm mở rộng.

\begin{table}[H]
\centering
\caption{Ảnh hưởng của tìm kiếm tương đồng mở rộng đối với TBM}
\begin{tabularx}{\textwidth}{Xrr}
\toprule
Phương pháp & TM trung bình tốt nhất trong năm & Chênh lệch\\
\midrule
Chỉ MMseqs2 & 0.2117 & --\\
MMseqs2 + tìm kiếm tổng hợp & \textbf{0.3072} & +0.0955\\
Ứng viên đầu theo phương pháp tái tạo, an toàn theo thời gian & 0.2983 & +0.0866 so với chỉ MMseqs2\\
\bottomrule
\end{tabularx}
\end{table}

Tìm kiếm mở rộng cải thiện 11 trong 12 đối tượng. R1116 là ngoại lệ duy nhất, giảm
0.0084. Trên R1108, TM-score tăng từ 0.3124 lên 0.4889; trên R1117v2, MMseqs2 không
tìm được khuôn mẫu hữu ích nhưng nhánh mở rộng tạo đủ năm cấu trúc và tăng điểm từ
0.0999 lên 0.4199. Vì các bước còn lại được giữ cố định, mức tăng trung bình 0.0955
được quy chủ yếu cho khả năng truy hồi khuôn mẫu tốt hơn. Kết quả cũng cho thấy việc
thêm một cơ chế tìm kiếm không bảo đảm cải thiện mọi RNA: khuôn mẫu xa vẫn có thể làm
xếp hạng nhiễu nếu giống trình tự nhưng không phản ánh đúng cấu trúc.

Để giải thích cơ chế thay vì chỉ báo cáo điểm cuối, một ablation mới được chạy trên 20
RNA hiệu chỉnh. MMseqs2 chỉ trả candidate ở 8/20 RNA; composite trả đủ ở 20/20. Trên
8 RNA có hit, MMseqs2 đạt best-of-five TM 0.678316, nhưng đây là trung bình có điều
kiện và không đại diện 12 RNA thất bại retrieval. Hợp MMseqs2--composite đạt 0.389705
trên đủ 20 RNA, so với 0.376060 của composite đơn lẻ.

\begin{table}[H]
\centering
\caption{Ablation nguồn truy hồi trên 20 RNA hiệu chỉnh}
\begin{tabular}{lrrrr}
\toprule
Nguồn & RNA có candidate & Availability & Best-5 khi có & Coverage\\
\midrule
MMseqs2 & 8/20 & 0.40 & 0.678316 & 0.930965\\
Composite & 20/20 & 1.00 & 0.376060 & 0.754893\\
Hợp hai nguồn & 20/20 & 1.00 & \textbf{0.389705} & 0.812411\\
\bottomrule
\end{tabular}
\end{table}

Component study không xác nhận rằng bộ trọng số composite là tối ưu. G-only đạt best-5
0.381177, còn full weighted đạt 0.376060; chênh $+0.005117$ theo hướng G-only có CI
$[-0.002178,+0.015004]$. Thay đổi trọng số $G,L$ quanh $\pm25\%$ chỉ làm điểm dao động
trong khoảng 0.375435--0.376500. Kết luận đúng mức là score composite ổn định và giúp
recall, không phải mọi component đều tăng TM. Khi cho phép template phát hành sau target,
best-5 tăng lạc quan 0.026961; đây là định lượng trực tiếp cho nhu cầu temporal filter.

\begin{table}[H]
\centering
\caption{Ablation reranking trên cùng pool composite high-recall}
\begin{tabular}{lrrrr}
\toprule
Ranking & Top-1 & Top-3 & Top-5 & PDB khác nhau\\
\midrule
Identity & 0.334450 & 0.363854 & 0.371900 & 4.85\\
Identity $\times$ coverage & \textbf{0.363720} & \textbf{0.378412} & \textbf{0.387558} & 4.95\\
$\times$ completeness & 0.363720 & 0.378412 & 0.387558 & 4.95\\
+ distinct-PDB & 0.363720 & 0.378412 & 0.387558 & 5.00\\
\bottomrule
\end{tabular}
\end{table}

Coverage là component reranking có hiệu ứng thực dụng lớn nhất ($\Delta$ best-5
$+0.015659$), dù CI còn chứa 0. Completeness không đổi ranking trên cohort có tọa độ
gần đầy đủ; distinct-PDB tăng đa dạng nguồn nhưng không tăng TM. Với gap-fill, heuristic
cong không khác tuyến tính ở gap 1--8 nucleotide, nhưng giảm SW-RMSD9 trung bình
0.036778~\angstrom{} ở gap 9--20, CI $[0.009289,0.067605]$. Với gap $>20$, phần lớn
là terminal nên hai cách dùng cùng extrapolation và cho kết quả bằng nhau.

\subsection{Đánh giá bên ngoài trên Kaggle}

Hai chương trình Kaggle được đóng băng trước khi xem điểm của tập riêng đạt:

\begin{table}[H]
\centering
\caption{Đánh giá bên ngoài của TBM và pipeline lai cuối}
\begin{tabular}{lrr}
\toprule
Phương pháp & Công khai & Riêng tư\\
\midrule
TBM & 0.60084 & 0.60175\\
3 TBM + 2 DRfold2 + Geometry v2 & \textbf{0.62809} & \textbf{0.61390}\\
Chênh lệch tuyệt đối & +0.02725 & +0.01215\\
\bottomrule
\end{tabular}
\end{table}

Phiên bản TBM gồm MMseqs2 và tìm kiếm mở rộng, lọc theo thời gian và PDB của chính
đối tượng, căn chỉnh lại, chuyển tọa độ, lấp khoảng trống, tinh chỉnh hình học 300 bước
và kiểm tra tệp CSV. Phiên bản này không chứa DRfold2 hay Boltz. Sổ tay tính
toán TBM dùng làm tham chiếu trên bảng xếp hạng có điểm được báo cáo là 0.59298, thấp
hơn 0.00877 so với điểm trên tập riêng tư 0.60175 của phương pháp trong luận văn. Do tập ẩn không cung
cấp điểm theo từng RNA, kết quả Kaggle là bằng chứng đánh giá bên ngoài cho toàn nhánh
TBM, không phải bằng chứng nhân quả riêng cho từng thành phần.

Phiên bản lai \texttt{rna3d-thesis-hybrid-geometry-v2}, phiên bản 10, giữ ba cấu trúc
TBM xếp hạng cao nhất và thay hai vị trí còn lại bằng hai cấu trúc DRfold2 có độ tin cậy
mô hình cao nhất khi chiều dài không quá 600 nucleotide; nếu DRfold2 không khả thi thì
vị trí đó quay về TBM. Geometry v2 sau đó tinh chỉnh cả năm cấu trúc trong 300 bước.
Trong lần chạy kiểm tra công khai, DRfold2 sinh đủ hai cấu trúc cho 11/12 RNA; RNA dài
720 nucleotide dùng năm cấu trúc TBM, và Geometry v2 hoàn tất cho 60/60 cấu trúc. Tệp
kết quả khớp đúng 2.515 ID mẫu, không trùng ID, không có tọa độ thiếu và không sử dụng
cấu trúc thực nghiệm. Đây là phép loại bỏ thành phần ở mức quy trình: mức tăng 0.01215
trên tập riêng chứng minh cấu hình lai có ích trên tập ẩn, nhưng không tách riêng phần
đóng góp của DRfold2 và Geometry v2.

\section{RQ2: Giá trị bổ sung của mô hình học sâu}

Sau khi tăng khả năng tìm khuôn mẫu, câu hỏi tiếp theo là liệu mô hình học sâu chỉ lặp
lại các nếp gấp đã có hay tạo thêm các giả thuyết hữu ích. Để tách chất lượng sinh cấu
trúc khỏi chất lượng lựa chọn, thí nghiệm báo cáo cả điểm lý tưởng và điểm được lựa chọn
không sử dụng cấu trúc thực nghiệm.

Bằng chứng chính đến từ 20 RNA mới nhất. Candidate tốt nhất trong ba TBM đạt TM
0.565183; candidate tốt nhất trong hai DRfold2 đạt 0.502627. DRfold2 vì vậy không mạnh
hơn TBM theo trung bình, nhưng thắng TBM trên 8/20 RNA. Khi đặt cả năm cấu trúc vào tập,
best-of-five/union oracle tăng lên 0.588304, chênh $+0.023121$ so với TBM với CI
$[+0.006444,+0.047519]$. Tám RNA cải thiện và 12 RNA giữ nguyên vì scorer vẫn có thể
chọn candidate TBM cũ.

\begin{table}[H]
\centering
\caption{Allocation TBM--DRfold2 với Geometry tắt trên 20 RNA mới nhất}
\begin{tabular}{lrrr}
\toprule
Allocation & Số candidate & Mean best TM & Chênh so với TBM cùng $N$\\
\midrule
1T+0D & 1 & 0.533635 & --\\
0T+1D & 1 & 0.488402 & $-0.045233$\\
2T+0D & 2 & 0.549041 & --\\
1T+1D & 2 & \textbf{0.579587} & $+0.030546$\\
3T+0D & 3 & 0.565183 & --\\
2T+1D & 3 & \textbf{0.581195} & $+0.016012$\\
1T+2D & 3 & 0.582752 & $+0.017569$\\
3T+2D & 5 & \textbf{0.588304} & $+0.023121$ so với 3T\\
\bottomrule
\end{tabular}
\end{table}

Fixed-$N=3$ thay một TBM bằng một DRfold2 tăng trung bình 0.016012 nhưng CI
$[-0.002072,+0.041466]$ chứa 0: hiệu ứng có lợi tập trung ở 7 RNA, 10 hòa và 3 giảm.
Ở $N=2$, 1T+1D tăng 0.030546 với CI $[+0.002949,+0.068185]$. Hai DRfold2 đơn lẻ lại
thấp hơn hai TBM 0.046414. Mẫu này hỗ trợ claim ``complementarity'': giữ đại diện của
cả hai nguồn tốt hơn thay toàn bộ TBM bằng deep learning.

\subsection{Bằng chứng phát triển CASP15}

\begin{table}[H]
\centering
\caption{Khả năng bao phủ nếp gấp của TBM và mô hình học sâu}
\small
\begin{tabularx}{\textwidth}{>{\raggedright\arraybackslash}X*{5}{>{\centering\arraybackslash}p{0.105\textwidth}}}
\toprule
Phân tích & Số RNA & TBM lựa chọn & Hợp lựa chọn & TBM lý tưởng & Hợp lý tưởng\\
\midrule
Đủ 12 RNA & 12 & 0.3113 & 0.4713 & 0.3143 & \textbf{0.5123}\\
Loại R1128 trùng trình tự & 11 & 0.3155 & 0.4245 & 0.3188 & \textbf{0.4693}\\
\bottomrule
\end{tabularx}
\end{table}

Mức tăng của điểm lý tưởng lần lượt là +0.1981 và +0.1505. Khi loại R1128 có trùng lặp trình tự,
mô hình học sâu vẫn làm tăng điểm này ở 10 trong 11 RNA và giữ nguyên một RNA. Do đó lợi
ích không chỉ đến từ trường hợp trùng dữ liệu. Tuy nhiên, nguồn học sâu không luôn vượt
TBM: ở R1117v2, điểm lý tưởng của TBM đạt 0.4616 còn mô hình học sâu chỉ đạt 0.2799. Hệ thống lai
vì thế cần bảo tồn cấu trúc tốt từ cả hai nguồn, thay vì luôn ưu tiên một nguồn.

Điểm lựa chọn của pilot đạt 0.4713, thấp hơn oracle 0.5123 khoảng 0.0411. Kết quả này
giải thích vì sao GeoFuse từng được khảo sát, nhưng không còn là phát hiện trung tâm của
luận văn sau khi phép hợp nhất không khái quát. Pipeline cuối giữ allocation cố định.

Độ tin cậy thô không giải quyết được nút thắt này. Trên 55 cấu trúc DRfold2, tương quan
Spearman gộp giữa độ tin cậy và TM là 0.486, nhưng tương quan trung bình khi xếp hạng
trong từng RNA chỉ là $-0.036$. Quan hệ gộp chủ yếu phân biệt những RNA dễ và khó, chứ
không ổn định để chọn giữa các cấu trúc của cùng một RNA hoặc giữa TBM với DRfold2.

\subsection{Boltz trên RNA dài}

R1138 dài 720 nucleotide:

\begin{equation}
\mathrm{TM}_{TBM}=0.2243,\quad
\mathrm{TM}_{Boltz}=0.2751,\quad
\Delta=+0.0508.
\end{equation}

Kết quả này cho thấy Boltz có thể cung cấp một phương án bổ sung khi DRfold2 không chạy
được vì giới hạn bộ nhớ. Một trường hợp đơn lẻ không đủ để kết luận Boltz luôn tốt hơn
DRfold2 hoặc TBM.

\iffalse
\section{RQ3: Tinh chỉnh độ chính xác hình học cục bộ}

Giả thuyết của Geometry v2 không phải tạo một nếp gấp mới mà là sửa sai lệch cục bộ
trong khi giữ cấu trúc nguồn. Vì các đại lượng như va chạm hoặc NLL đồng thời là thành
phần của hàm mục tiêu, chúng không được dùng làm bằng chứng chính. Thay vào đó, cùng 60
cấu trúc của 12 RNA được so sánh trước và sau tinh chỉnh bằng TM-score,
\cprime{}-lDDT và RMSD cục bộ so với cấu trúc thực nghiệm.

\begin{table}[H]
\centering
\caption{Đánh giá Geometry v2 bằng các chỉ số độc lập}
\begin{tabular}{lrrr}
\toprule
Chỉ số & Trước tinh chỉnh & Geometry v2 & Mức thay đổi theo chiều tốt\\
\midrule
TM tốt nhất trong năm & 0.471268 & 0.471308 & +0.000040\\
\cprime{} RMSD (\angstrom) & 28.2718 & 28.2650 & +0.006893\\
\cprime{}-lDDT & 0.472117 & \textbf{0.481823} & \textbf{+0.009706}\\
RMSD cửa sổ 9 (\angstrom) & 4.00131 & \textbf{3.94415} & \textbf{+0.057166}\\
RMSD cửa sổ 15 (\angstrom) & 6.54814 & 6.52453 & +0.023608\\
RMSD cửa sổ 31 (\angstrom) & 11.8667 & 11.8716 & $-0.004926$\\
\bottomrule
\end{tabular}
\end{table}

\cprime{}-lDDT tăng trên cả 12 RNA và 56 trong 60 cấu trúc; khoảng tin cậy 95\% của
chênh lệch là $[+0.007403,+0.012617]$. RMSD cửa sổ 9 cũng cải thiện trên cả 12 RNA,
với khoảng tin cậy $[+0.035077,+0.091680]$~\angstrom{} theo chiều giảm lỗi. Ngược lại,
chênh lệch TM có khoảng tin cậy $[-0.003888,+0.003962]$ và RMSD cửa sổ 31 gần như
không đổi. Mẫu kết quả này phù hợp với mục tiêu thiết kế: Geometry v2 điều chỉnh mạng
khoảng cách ở quy mô ngắn nhưng không tạo ra thay đổi đáng kể ở bố cục toàn cục. Do đó
đóng góp của phương pháp nằm ở tinh chỉnh cục bộ, không phải cải thiện nếp gấp tổng thể.

\fi

\section{RQ3: Hiệu quả, cơ chế và khả năng khái quát của Geometry v2}

\subsection{Khóa cấu hình trên 20 RNA hiệu chỉnh}

Priors được ước lượng lại chỉ từ 60 RNA train. Trên 20 RNA calibration, hai cấu hình
full fixed/adaptive đều qua gate: TM không giảm quá 0.005, lDDT tăng và ít nhất hai
trong ba sliding-window RMSD giảm. Source restraint 1.5 tạo mức tăng lDDT lớn nhất trong
các cấu hình full khai báo trước nên được đóng băng trước khi mở validation.

\begin{table}[H]
\centering
\caption{Các ablation tiêu biểu trên 20 RNA hiệu chỉnh}
\small
\begin{tabular}{lrrrrrr}
\toprule
Variant & Best-5 TM & lDDT & SW9 & Kink & Angle NLL & Torsion NLL\\
\midrule
Raw & 0.407436 & 0.500573 & 4.175314 & 0.043392 & 0.867374 & 0.790828\\
Source+backbone & 0.406346 & 0.510639 & 4.170749 & 0.063631 & 0.910474 & 0.823790\\
Full fixed & 0.404957 & 0.506502 & 4.148670 & 0.039937 & 0.630507 & 0.704167\\
Full adaptive & 0.404908 & 0.508102 & 4.151196 & 0.040512 & 0.635805 & 0.702107\\
Full, source=1.5 & 0.404674 & \textbf{0.511698} & \textbf{4.144572} & \textbf{0.038782} & \textbf{0.591498} & \textbf{0.677023}\\
\bottomrule
\end{tabular}
\end{table}

Đối chứng source+backbone thậm chí tăng lDDT mạnh, nhưng làm sharp-kink, angle NLL và
torsion NLL xấu hơn. Leave-one-out giải thích các thành phần: bỏ backbone đưa lDDT về
0.500739, gần raw; bỏ angle làm angle NLL tăng từ 0.638915 lên 0.837255; bỏ torsion làm
torsion NLL tăng từ 0.706165 lên 0.829415; bỏ kink guard tăng kink từ 0.040512 lên
0.045454; bỏ clash/Rg tăng clash proxy. Không component nào một mình chứng minh native
accuracy; vai trò của chúng là kiểm soát kiểu lỗi tương ứng.

Convergence cho thấy lDDT calibration tăng lần lượt 0.506190, 0.506357, 0.506953,
0.507858 và 0.508102 ở 25, 50, 100, 200 và 300 bước. Mức cải thiện dần bão hòa, không
có dấu hiệu cần tối ưu vô hạn. Learning rate 0.02--0.08 và source restraint 1.5--6.0
đều giữ TM trong gate; phương pháp không phụ thuộc đúng một hằng số duy nhất.

\subsection{Kết quả confirmatory trên 20 RNA mới nhất}

\begin{table}[H]
\centering
\caption{Raw, đối chứng đơn giản và Geometry v2 trên tập xác nhận mới nhất}
\small
\begin{tabular}{lrrr}
\toprule
Chỉ số & Raw & Source+backbone & Full Geometry v2\\
\midrule
Best-of-5 TM & 0.588304 & 0.587767 & 0.587197\\
\cprime{} RMSD (\angstrom) & 10.903641 & 10.902999 & 10.903764\\
\cprime{}-lDDT & 0.659836 & \textbf{0.667417} & 0.666496\\
SW-RMSD9 (\angstrom) & 2.651395 & 2.641594 & \textbf{2.611010}\\
SW-RMSD15 (\angstrom) & 4.273393 & 4.264781 & \textbf{4.256136}\\
SW-RMSD31 (\angstrom) & 7.562650 & 7.558546 & \textbf{7.557402}\\
Sharp-kink fraction & 0.048977 & 0.068558 & \textbf{0.044498}\\
Angle NLL & 0.820552 & 0.869599 & \textbf{0.618824}\\
Torsion NLL & 0.837811 & 0.871896 & \textbf{0.734876}\\
\bottomrule
\end{tabular}
\end{table}

Full Geometry v2 tăng lDDT $+0.006660$, CI 95\%
$[+0.000933,+0.013093]$, cải thiện 12 và làm giảm 8 RNA. SW-RMSD9 giảm 0.040386,
CI $[0.029354,0.052090]$, cải thiện 19/20 RNA; SW-RMSD15 giảm 0.017257,
CI $[0.000617,0.031336]$, cải thiện 15/20. SW-RMSD31 giảm 0.005248 nhưng CI chứa 0.
TM giảm 0.001107, CI $[-0.002809,+0.000418]$: không có bằng chứng TM thay đổi và mức
giảm nằm trong biên bảo toàn 0.005. Gate confirmatory vì vậy đạt.

Đối chứng đơn giản tăng lDDT $+0.007581$, lớn hơn full 0.000921 nhưng chênh trực tiếp có
CI $[-0.003118,+0.001327]$; full không thắng đối chứng về lDDT. Full thắng rõ ở SW9
($+0.030585$ theo chiều giảm lỗi, CI $[0.019640,0.044115]$) và sửa các diagnostic mà
đối chứng làm xấu. Claim đúng là full Geometry tạo projection cân bằng hơn ở nhiều
khía cạnh; không phải mọi metric native đều vượt baseline đơn giản.

\subsection{Factorial nguồn ứng viên và Geometry}

Factorial fixed-$N=3$ cho kết quả: TBM 3T+0D đạt 0.565183 khi Geometry off và 0.559331
khi on; hybrid 2T+1D đạt 0.581195 off và 0.578592 on. Hiệu ứng Geometry là
$-0.005852$ trên TBM (CI chứa 0) và $-0.002604$ trên hybrid (CI
$[-0.005473,-0.000278]$). Interaction $+0.003248$ có CI chứa 0. Như vậy mức tăng của
pipeline Kaggle không được quy cho Geometry: nguồn candidate đa dạng tăng fold coverage,
còn Geometry chủ yếu sửa cục bộ và có thể đổi candidate thắng best-of-$N$ theo chiều xấu.

\subsection{Khi nào tinh chỉnh có lợi?}

Theo chất lượng raw, nhóm thấp tăng lDDT $+0.016261$, CI
$[+0.008900,+0.024089]$; nhóm cao thay đổi $-0.001788$, CI chứa 0. Nhóm TBM tăng
$+0.007610$ với CI dương; DRfold2 tăng $+0.005235$ nhưng CI chứa 0. Hiệu ứng rõ hơn ở
candidate support thấp và confidence trung bình; RNA ngắn/dài có interval rộng do chỉ
có 4/5 target. Phát hiện này gợi ý Geometry nên được gate theo chất lượng ban đầu trong
tương lai, nhưng luận văn không học gate mới sau khi đã mở validation.

\iffalse
\section{RQ4: Từ lựa chọn nguồn đến hợp nhất tọa độ}

RQ4 gồm hai kiểm định tách biệt. Kiểm định thứ nhất hỏi mô hình có nhận ra nguồn tốt hơn
theo nucleotide hay không. Kiểm định thứ hai hỏi liệu dự đoán đó có thể tạo một cấu trúc
hợp nhất tốt hơn hay không. Việc tách hai bước tránh suy luận rằng độ chính xác phân loại
tự động đồng nghĩa với độ chính xác cấu trúc.

\subsection{Lựa chọn đại lượng giám sát}

Trên 100 RNA, 352 cặp nguồn và 24,876 nucleotide hợp lệ, quyết định theo sai lệch điểm
sau căn chỉnh đồng thuận với \cprime{}-lDDT ở 69.7\% vị trí; tương quan Spearman sau
khi loại ảnh hưởng giữa các RNA là 0.319. Mức đồng thuận giữa sai lệch điểm và RMSD cửa
sổ 15 là 72.7\%, còn giữa lDDT và RMSD cửa sổ là 70.9\%. Các đại lượng chia sẻ một phần
tín hiệu nhưng không tương đương. Vì \cprime{}-lDDT đo quan hệ khoảng cách cục bộ mà
không phụ thuộc một phép chồng khít toàn cấu trúc, nó được chọn làm đại lượng giám sát
chính.

\subsection{Khả năng lựa chọn nguồn trên RNA mới}

\begin{table}[H]
\centering
\caption{Ước lượng nguồn tốt hơn trên 20 RNA mới nhất}
\small
\begin{tabularx}{\textwidth}{>{\raggedright\arraybackslash}X>{\centering\arraybackslash}p{0.23\textwidth}>{\centering\arraybackslash}p{0.14\textwidth}>{\centering\arraybackslash}p{0.12\textwidth}}
\toprule
Phương pháp & Lỗi $1-$lDDT trung bình theo RNA & Độ chính xác & ROC-AUC\\
\midrule
Lựa chọn lý tưởng theo nucleotide & 0.193424 & 0.923691 & 0.934561\\
Conv1D & \textbf{0.238994} & 0.650573 & 0.729868\\
Hồi quy logistic & 0.247937 & 0.628682 & 0.687387\\
Tăng cường gradient & 0.250951 & 0.643208 & 0.731747\\
Luôn chọn DRfold2 & 0.269575 & 0.583061 & 0.500000\\
Luật khoảng trống & 0.322453 & 0.458265 & 0.534949\\
Luật độ tin cậy & 0.326173 & 0.452332 & 0.530001\\
Luôn chọn TBM & 0.347735 & 0.416939 & 0.500000\\
\bottomrule
\end{tabularx}
\end{table}

Tập hiệu chỉnh chọn Conv1D với ngưỡng 0.75. So với luật dựa trên khoảng trống được chọn
trên cùng tập hiệu chỉnh, lỗi trung bình giảm 0.083458; khoảng tin cậy 95\% kéo dài từ
+0.023419 đến +0.155562, với 14 trong 20 RNA cải thiện. Như vậy mô hình học được một
tín hiệu có khả năng khái quát để so sánh chất lượng cục bộ của hai nguồn. Tuy nhiên,
AUC 0.729868 và độ chính xác 0.650573 cũng cho thấy dự đoán còn xa mức hoàn hảo, nên cần
một cơ chế kiểm soát rủi ro trước khi sửa tọa độ.

\subsection{Kết quả hợp nhất có chọn lọc}

Ba ngưỡng self-TM 0.35, 0.45 và 0.55 cho cùng kết quả tổng hợp trên tập hiệu chỉnh;
0.35 được chọn theo thứ tự định trước. Cụm chứa cả TBM và DRfold2 xuất hiện ở 3/20 RNA
hiệu chỉnh và 12/20 RNA xác nhận. Do đó thiếu cặp cùng nếp gấp không phải nguyên nhân
duy nhất khiến hợp nhất khó xảy ra.

\begin{table}[H]
\centering
\caption{Kết quả hợp nhất có chọn lọc trên 20 RNA mới nhất}
\small
\begin{tabularx}{\textwidth}{>{\raggedright\arraybackslash}X*{4}{>{\centering\arraybackslash}p{0.14\textwidth}}}
\toprule
Biến thể & TM được chọn & lDDT được chọn & TM lý tưởng & lDDT lý tưởng\\
\midrule
F0 cấu trúc gốc & 0.588304 & 0.795059 & 0.588304 & 0.795059\\
F1 quy tắc kinh nghiệm & 0.589984 & 0.793841 & 0.592826 & 0.795828\\
F2 lựa chọn theo chất lượng & 0.588304 & 0.795059 & 0.588304 & 0.795059\\
F3 F2 + Geometry v2 & 0.588304 & 0.795059 & 0.588304 & 0.795059\\
F4 dẫn hướng bằng cấu trúc thực nghiệm & 0.589158 & 0.795648 & 0.589160 & 0.795648\\
\bottomrule
\end{tabularx}
\end{table}

F2 và F3 từ chối mọi cặp vì không có đoạn nào đồng thời vượt xác suất 0.90, nằm trong
giới hạn bất đồng và kéo dài ít nhất 7 nucleotide. Vì vậy cả 20 RNA giữ nguyên kết quả
của F0. Cơ chế bảo thủ đã tránh suy giảm, nhưng không cung cấp bằng chứng rằng phép hợp
nhất cải thiện cấu trúc.

F1 tạo 24 cấu trúc trên 12 RNA và tăng TM lựa chọn trung bình 0.001681, nhưng khoảng
tin cậy $[-0.004908,+0.008495]$ chứa 0, trong khi lDDT lựa chọn giảm 0.001218. Kết quả
nhỏ và không ổn định này không đủ để thay thế phương pháp bảo thủ bằng quy tắc kinh
nghiệm. Dữ liệu dẫn tới kết luận:

Kết quả này làm rõ phát hiện trung tâm của RQ4: mô hình có thể nhận ra nguồn tốt hơn
theo \cprime{}-lDDT cục bộ, nhưng khả năng đó chưa chuyển thành một phép hợp nhất tọa
độ hữu ích. Biết nguồn nào tốt hơn không đồng nghĩa biết cắt, căn chỉnh và ghép hai cấu
trúc sao cho toàn cấu trúc tốt hơn.

\subsection{Minh họa trên R1108}

Cấu trúc TBM tốt nhất của R1108 là 1DRZ\_B, TM 0.49384; DRfold2 P1 đạt 0.60192, tăng
0.10808 khi đặt hai nguồn vào cùng tập. Geometry v2 đưa P1 từ TM 0.60192 lên 0.60227
và lDDT từ 0.788065 lên 0.788827. Trường hợp này minh họa đúng vai trò khác nhau của
hai thành phần: DRfold2 cung cấp một nếp gấp tốt hơn khuôn mẫu, còn Geometry v2 chỉ điều
chỉnh cục bộ trên nếp gấp đó.

\subsection{Minh họa trên 8K0Y\_A}

RNA dài 56 nucleotide có ba cấu trúc TBM và hai cấu trúc DRfold2. Cặp 8V1I\_B--P1 có
RMSD sau căn chỉnh 1.099~\angstrom, độ bất đồng trung bình 2.855~\angstrom{} và xác suất
trung bình chọn DRfold2 là 0.7218. Không có đoạn dài 7 nucleotide vượt 0.90, nên F2
không hợp nhất. TM/lDDT của tập gốc là 0.65667/0.871251 và được giữ nguyên. Đây là ví
dụ về sự đánh đổi: quy tắc an toàn bảo vệ dự đoán tốt, nhưng đồng thời không tạo giá trị
mới khi tín hiệu chưa đủ mạnh.

\fi

\section{Mạch phát hiện qua ba câu hỏi nghiên cứu}

Ba câu hỏi tạo thành một chuỗi rõ: composite giúp TBM không rơi vào trạng thái không có
khuôn; DRfold2 bổ sung một số nếp gấp TBM bỏ lỡ; Geometry v2 sửa quan hệ cục bộ trên
candidate đã có. Candidate diversity cải thiện best-of-five TM, còn Geometry không tạo
nếp gấp mới và không tăng TM. Giá trị khoa học nằm ở việc tách đúng hai hiệu ứng này và
kiểm chứng local refinement trên một split thời gian/họ độc lập.

% -----------------------------------------------------------------------------
% CHAPTER 7
% -----------------------------------------------------------------------------
\chapter{Thảo luận}

\section{Từ retrieval coverage đến candidate complementarity}

RQ1 cho thấy hai tầng search có vai trò khác nhau. MMseqs2 tạo candidate chất lượng cao
khi có hit nhưng chỉ phủ 40\% cohort hiệu chỉnh; composite hạ precision trung bình để
đạt availability 100\%. Reranking identity$\times$coverage chuyển một phần recall đó
thành top-$k$ precision. Distinct-PDB không tăng TM trên cohort này nhưng làm năm giả
thuyết ít phụ thuộc một PDB hơn. Đây là kết luận component-specific, mạnh hơn việc chỉ
nói ``full TBM tăng điểm''.

RQ2 cho thấy DRfold2 không thay thế TBM: oracle riêng của DRfold2 thấp hơn, nhưng thắng
ở 8/20 RNA. Giữ đại diện cả hai nguồn làm union tăng 0.023121. Fixed-size allocation
xác nhận hướng hiệu ứng dương nhưng CI còn rộng; do đó claim là complementarity có điều
kiện theo target, không phải deep learning luôn ưu việt.

\section{Vì sao Geometry v2 cải thiện cục bộ nhưng không tăng TM-score}

Hàm mục tiêu của Geometry v2 được thiết kế để giữ gần cấu trúc nguồn. Trên tập
confirmatory, lDDT và RMSD cửa sổ 9/15 cải thiện với CI không chứa 0, trong khi TM giảm
0.001107 và CI chứa 0. Projection dịch mỗi nucleotide trung bình khoảng
0.305~\angstrom{} nhưng không tạo tiếp xúc xa hay nếp gấp mới.

TM-score ưu tiên nếp gấp tổng thể và giảm ảnh hưởng của một số sai số cục bộ. Geometry
v2 không tạo tiếp xúc xa mới, không đổi cách các miền sắp xếp và không tìm một kiểu nếp gấp
khác; do đó không có cơ chế mạnh để tăng TM-score nếu cấu trúc nguồn sai toàn cục. Giá
trị của phương pháp nằm ở việc sửa vừa phải quan hệ \cprime{} cục bộ mà không làm mất
giả thuyết nếp gấp đã có. Đối chứng source--backbone cho thấy một phần lDDT gain đến từ
smoothing đơn giản; full Geometry được biện minh bởi SW9 tốt hơn và việc không làm xấu
kink/angle/torsion. Stratification cho thấy gain tập trung ở candidate yếu, còn nhóm
chất lượng cao không cải thiện. Vì biểu diễn chỉ gồm \cprime{}, kết quả không cho phép
suy rộng sang bắt cặp base hay hóa học lập thể toàn nguyên tử.

\iffalse
\section{Vì sao lựa chọn nguồn chưa chuyển thành hợp nhất tọa độ}

Kết quả Conv1D cho thấy các đặc trưng độ tin cậy, hỗ trợ khuôn mẫu, bất đồng và hình học
có chứa thông tin về nguồn tốt hơn theo nucleotide. Tuy nhiên, nhãn phân loại chỉ trả
lời một câu hỏi so sánh: tại vị trí này, nguồn nào có quan hệ cục bộ gần cấu trúc thực
nghiệm hơn? Nó không trực tiếp cho biết vector dịch chuyển nào nên áp dụng, biên đoạn
nên đặt ở đâu hoặc cách bảo toàn tương quan giữa các miền.

Phép hợp nhất tuyến tính còn phụ thuộc vào chồng khít hai cấu trúc. Khi hai miền có
chuyển động tương đối, một phép Kabsch toàn cục có thể căn miền này tốt nhưng để miền
khác lệch. Một đoạn có lDDT tốt hơn trong DRfold2 cũng không nhất thiết khớp tọa độ và
hướng tiếp tuyến với hai đầu của đoạn TBM. Lấy trung bình có thể tạo một cấu trúc nằm
giữa hai giả thuyết chứ không gần cấu trúc đúng. Vì vậy, lựa chọn chất lượng và kiến tạo
tọa độ là hai bài toán khác nhau.

Ngưỡng bảo thủ của F2 đã thể hiện đúng vai trò kiểm soát rủi ro nhưng quá nghiêm để tạo
dữ liệu về hiệu quả: không quyết định nào được thực hiện. Nới ngưỡng sau khi nhìn tập
xác nhận có thể tạo một mức tăng nhỏ, nhưng sẽ làm mất tính độc lập của đánh giá. Luận
văn vì thế giữ nguyên kết quả âm và xem nó như chỉ dẫn cho một phương pháp hợp nhất học
phần dư tọa độ hoặc dùng hệ tọa độ cục bộ trong nghiên cứu tiếp theo.

\fi
\section{Quan hệ giữa phần kế thừa và đóng góp của luận văn}

\begin{table}[H]
\centering
\caption{Nguồn gốc và vai trò của các thành phần}
\begin{tabularx}{\textwidth}{p{0.25\textwidth}p{0.27\textwidth}X}
\toprule
Thành phần & Nguồn ý tưởng & Phần được thực hiện hoặc nghiên cứu thêm\\
\midrule
Tìm kiếm tương đồng mở rộng & Sổ tay tính toán TBM tham chiếu & Kết hợp MMseqs2, lọc thời gian/PDB, xếp hạng lại và đánh giá có đối chứng.\\
Chuyển tọa độ và lấp khoảng trống & TBM truyền thống và mã cuộc thi & Mặt nạ hỗ trợ, độ tin cậy theo nucleotide và biểu diễn nguồn gốc thống nhất.\\
DRfold2/Boltz & Mô hình công bố & Kiểm tra mốc dữ liệu/trùng lặp, chuẩn hóa ứng viên và allocation fixed-$N$.\\
Geometry v2 & Đóng góp phương pháp & Phân bố tiên nghiệm theo ngữ cảnh, cường độ thích nghi, ràng buộc nguồn, bảo vệ góc gập và đánh giá cục bộ độc lập.\\
Quy trình đánh giá & Các bộ đánh giá chuẩn về cấu trúc & Chia theo thời gian/họ, nhiều chỉ số bổ sung và lấy mẫu lại theo RNA.\\
\bottomrule
\end{tabularx}
\end{table}

\section{Quan hệ với kết quả cuộc thi}

Sổ tay tính toán TBM tham chiếu có điểm được báo cáo là 0.59298. Phiên bản lai V4 có điểm
0.57631 và kết hợp Boltz, DRfold2 với ràng buộc cùng phương án quay về khuôn mẫu, nhưng
không thực hiện lựa chọn nguồn cục bộ theo độ tin cậy. Mã nguồn này cũng có một số vấn đề
về chỉ số hàng, lựa chọn đối tượng và tọa độ giữ chỗ; vì vậy điểm thấp hơn không phải
bằng chứng rằng mô hình học sâu về bản chất kém hơn TBM
\cite{top1_tbm,top1_hybrid}.

Phương pháp TBM của luận văn đạt 0.60175 trên tập riêng tư. Pipeline lai cuối
đạt 0.61390, cao hơn 0.01215, nên khả năng triển khai và lợi ích của toàn cấu hình lai đã
được xác nhận bằng một lần nộp bên ngoài. Tuy nhiên, kết quả này đồng thời thay đổi thành
phần ứng viên và áp dụng Geometry v2; vì vậy không được dùng để quy toàn bộ mức tăng cho
một thành phần riêng. Factorial công khai cho thấy source diversity làm tăng fold
coverage còn Geometry không tăng TM; bằng chứng cho Geometry đến từ chỉ số cục bộ trên
tập xác nhận, không phải từ việc dò ngược bảng xếp hạng.

\section{Các nguy cơ đối với tính hợp lệ}

\subsection{Tính hợp lệ nội tại}

Nguy cơ gồm rò rỉ qua ngày công bố, khuôn mẫu của chính đối tượng, trình tự trùng hoàn
toàn, mốc huấn luyện của mô hình học sâu và việc thay đối tượng sau lỗi kỹ thuật. Các
kiểm tra đã giảm rủi ro nhưng không thể chứng minh tuyệt đối rằng bộ trọng số chưa học
một quan hệ họ hàng xa. Một số dữ liệu lớn nằm ngoài Git nên cần tệp kê khai và mã kiểm tra
đầy đủ trước khi công bố.

\subsection{Khả năng khái quát}

Geometry v2 được xác nhận trên 20 RNA dài tối đa 96 nucleotide, sau calibration 20 RNA
và prior-train 60 RNA. Kết quả chưa chắc áp dụng cho RNA rất dài, phức hợp nhiều chuỗi
hoặc họ RNA ít được đại diện. Tập ẩn Kaggle khác CASP15 nên điểm giữa hai tập không thể
so sánh trực tiếp.

\subsection{Mức độ phù hợp của đại lượng đo}

\cprime{}-lDDT là bản thích nghi, độ dài cửa sổ là lựa chọn phân tích, còn tỷ lệ góc gập
sắc, va chạm \cprime{}, NLL góc phẳng/góc xoắn và chỉ dấu có dạng cặp là các đại lượng
nội bộ. Chúng không đo trực tiếp độ đúng của cặp base nitơ, xếp chồng hoặc tính hợp lý vật lý toàn nguyên tử.

\subsection{Tính hợp lệ của kết luận thống kê}

Số RNA độc lập nhỏ hơn số nucleotide rất nhiều; candidate của cùng RNA cũng tương quan.
Mỗi subgroup trước hết được trung bình trong target/group cell rồi mới bootstrap RNA.
Những CI chứa 0 được trình bày là bằng chứng chưa đủ, không bị diễn giải thành cải thiện.

\section{Hạn chế kỹ thuật}

DRfold2 đòi hỏi GPU và khó chạy trên trình tự dài; Boltz chưa được tích hợp thống nhất
vào đường xác nhận. Bản Kaggle lai vì vậy giới hạn DRfold2 ở 600 nucleotide và quay về
TBM cho trình tự dài hơn. Bộ chọn năm cấu trúc vẫn dùng một luật cố định ba TBM cộng hai
DRfold2 thay vì một ước lượng chất lượng toàn cấu trúc đã hiệu chỉnh.
Một số nucleotide biến đổi hiếm được ánh xạ thành N, có thể làm giảm nhẹ mức đồng nhất
trình tự dù tọa độ vẫn được chuyển.

Trong phạm vi những giới hạn trên, kết luận của luận văn được giữ ở mức tương xứng với
dữ liệu: TBM và pipeline lai cuối đều đã được đánh giá bên ngoài trên Kaggle; nguồn học
sâu tăng fold coverage; Geometry v2 cải thiện đường \cprime{} cục bộ trên tập
confirmatory nhưng không tăng TM-score.

% -----------------------------------------------------------------------------
% CHAPTER 8
% -----------------------------------------------------------------------------
\chapter{Kết luận và hướng phát triển}

\section{Kết luận}

Luận văn đã xây dựng một pipeline RNA 3D mức \cprime{} an toàn theo thời gian. MMseqs2
tạo precision khi có homolog rõ, còn composite bảo đảm retrieval coverage; reranking
coverage chuyển một phần recall thành chất lượng top-$k$. Nhánh TBM đạt private Kaggle
0.60175. DRfold2 không mạnh hơn TBM trên mọi RNA nhưng thắng ở 8/20 RNA mới nhất; union
3T+2D đạt oracle 0.588304 so với 0.565183 của 3T. Pipeline lai cuối đạt private Kaggle
0.61390. Các kết quả fixed-$N$ cho phép quy lợi ích này cho candidate complementarity
thận trọng hơn, thay vì gán toàn bộ cho deep learning hoặc Geometry.

Đóng góp phương pháp Geometry v2 được kiểm tra bằng protocol 60/20/20. Trên 20 RNA
confirmatory, lDDT tăng 0.006660 và SW-RMSD9 giảm 0.040386~\angstrom{} với CI dương;
TM giảm 0.001107, CI chứa 0 và nằm trong biên bảo toàn. Đối chứng source--backbone cho
lDDT tương đương, nhưng full Geometry tốt hơn rõ ở SW9 và tránh làm xấu kink/angle/
torsion. Hiệu ứng tập trung ở candidate ban đầu yếu. Factorial không cho thấy Geometry
tăng TM, nên claim cuối là local refinement có kiểm chứng, không phải fold predictor.

Câu chuyện luận văn vì vậy gọn thành ba contribution: TBM time-safe có component
analysis; bằng chứng TBM--DRfold2 complementarity; và Geometry v2 với ablation,
baseline và held-out validation. GeoFuse bị loại khỏi core vì chưa móc được quality
signal thành một cải thiện tọa độ khái quát.

\section{Hướng phát triển ưu tiên}

\begin{enumerate}
  \item \textbf{Gate Geometry theo chất lượng ban đầu:} kiểm tra trên một cohort mới
  quy tắc chỉ refine candidate quality thấp; không học gate từ 20 RNA đã mở.
  \item \textbf{Mở rộng allocation bank:} sinh thật ít nhất năm TBM và năm DRfold2 mỗi
  RNA để chạy fixed-five 5T--0D tới 0T--5D, thay vì đệm bản sao.
  \item \textbf{Mở rộng số họ và chiều dài:} xác nhận trên RNA dài, phức hợp nhiều chuỗi
  và các họ ít đại diện; đánh giá riêng fallback Boltz.
  \item \textbf{Đánh giá toàn nguyên tử:} dựng cấu trúc đầy đủ để đo cặp base nitơ, xếp chồng,
  va chạm lập thể và độ trung thành của mạng tương tác khi tài nguyên cho phép.
  \item \textbf{Cải thiện lựa chọn cấu trúc:} học điểm chất lượng toàn cấu trúc có hiệu
  chỉnh riêng theo nguồn và mục tiêu đa dạng tường minh để giảm khoảng cách tới trần lý tưởng.
  \item \textbf{Đối chứng refinement ngoài hệ:} nếu dựng được full-atom, so sánh với
  QRNAS hoặc force-field refinement theo cùng source restraint và native-blind protocol.
\end{enumerate}

% -----------------------------------------------------------------------------
% APPENDICES
% -----------------------------------------------------------------------------
\appendix

\chapter{Thí nghiệm khám phá GeoFuse (không thuộc pipeline cuối)}

\section{Giả thuyết}

GeoFuse thử trả lời câu hỏi: nếu một đoạn TBM tốt còn đoạn khác DRfold2 tốt, liệu có thể
dùng đặc trưng native-blind để chọn nguồn theo nucleotide rồi ghép hai cấu trúc hay
không? Hai cấu trúc trước hết được phân cụm self-TM để tránh ghép khác nếp gấp, chồng
khít Kabsch, sau đó Conv1D dự đoán xác suất DRfold2 tốt hơn theo \cprime{}-lDDT cục bộ.
Chỉ đoạn dài ít nhất bảy nucleotide, xác suất ít nhất 0.90 và bất đồng 0.5--12~\angstrom{}
mới được ghép; nếu không, phương pháp từ chối.

\section{Kết quả khám phá}

Trên 20 RNA validation, Conv1D đạt ROC-AUC 0.729868 và lỗi $1-$lDDT 0.238994, tốt hơn
luật confidence/gap cố định. Tuy nhiên, điều đó chỉ chứng minh đặc trưng có tín hiệu so
sánh nguồn, chưa chứng minh vector tọa độ hợp nhất đúng. Với ngưỡng khóa trên
calibration, GeoFuse từ chối toàn bộ quyết định; TM/lDDT giữ nguyên 0.588304/0.795059.
Một luật kinh nghiệm nới hơn tạo 24 cấu trúc, tăng TM trung bình 0.001681 nhưng CI
$[-0.004908,+0.008495]$ chứa 0 và lDDT giảm 0.001218.

\section{Quyết định phạm vi}

GeoFuse không đạt tiêu chí ``tạo cấu trúc tốt hơn trên RNA mới'' và không tham gia
submission Kaggle 0.61390. Vì vậy luận văn không coi Conv1D, clustering hay fusion là
contribution. Thí nghiệm được lưu để minh bạch lịch sử nghiên cứu và giải thích một
failure mode: biết nguồn nào có lỗi cục bộ thấp hơn không tự động cho biết cách cắt,
căn chỉnh và ghép tọa độ liên tục. Hướng này chỉ nên mở lại với target mới, gate mới
được preregister và một mô hình học phần dư tọa độ SE(3)-equivariant.

\chapter{Bảng thuật ngữ}

\begin{longtable}{p{0.25\textwidth}p{0.65\textwidth}}
\toprule
Thuật ngữ & Nghĩa trong luận văn\\
\midrule
\endhead
Cấu trúc ứng viên & Một giả thuyết cấu trúc ba chiều cho RNA cần dự đoán.\\
Cấu trúc thực nghiệm & Cấu trúc được đo và dùng làm tham chiếu khi huấn luyện hoặc đánh giá.\\
Khuôn mẫu & Cấu trúc RNA đã biết dùng để chuyển tọa độ sang đối tượng mới.\\
TBM & Template-based modeling, mô hình hóa dựa trên khuôn mẫu.\\
Mô hình tiền huấn luyện & Mô hình đã học từ dữ liệu trước đó; bộ trọng số được giữ cố định khi suy luận.\\
Nếp gấp tổng thể & Kiểu nếp gấp và bố cục lớn của cấu trúc ba chiều.\\
Hình học cục bộ & Khoảng cách, góc phẳng và góc xoắn trong một vùng nucleotide ngắn.\\
Căn chỉnh & Ánh xạ nucleotide của đối tượng với nucleotide của khuôn mẫu.\\
Khoảng trống & Vị trí không có nucleotide khuôn mẫu tương ứng.\\
$k$-mer & Đoạn trình tự liên tiếp dài $k$.\\
Tìm kiếm mở rộng & Quét toàn thư viện tương đồng, không phải vét cạn không gian 3D.\\
Độ tin cậy & Tín hiệu nội bộ của mô hình, không đồng nhất với độ chính xác.\\
Hiệu chỉnh & Liên hệ đầu ra độ tin cậy với tần suất đúng quan sát được.\\
Oracle & Điểm lý tưởng do lựa chọn bằng cấu trúc thực nghiệm sau đánh giá; không triển khai được khi suy luận.\\
OOF & Dự đoán cho một mẫu nằm ngoài nếp gấp dữ liệu dùng để huấn luyện mô hình đó.\\
An toàn theo thời gian & Chỉ dùng dữ liệu đã tồn tại trước mốc của đối tượng.\\
Self-TM & TM-score giữa hai cấu trúc dự đoán, không cần cấu trúc thực nghiệm.\\
Cơ chế từ chối & Không hợp nhất khi bằng chứng chưa đủ mạnh.\\
Tinh chỉnh & Sửa nhẹ một cấu trúc có sẵn, không dự đoán lại nếp gấp từ đầu.\\
\bottomrule
\end{longtable}

\chapter{Danh sách kiểm tra khả năng tái lập}

\begin{enumerate}
  \item Tạo môi trường conda từ \texttt{environment.yml} và ghi phiên bản gói chính xác.
  \item Kiểm tra mã kiểm tra của CSV, CIF, cơ sở khuôn mẫu, bộ trọng số và dữ liệu tạm.
  \item Xây cơ sở khuôn mẫu bằng bộ đọc hiện tại; xác nhận số dòng và độ phủ ngày công bố.
  \item Kiểm tra rò rỉ thời gian và PDB của chính đối tượng; mốc thiếu hoặc sai phải bị từ chối.
  \item Chạy kiểm thử cho bộ đọc, căn chỉnh, chỉ số, gap-fill baseline và Geometry v2.
  \item Tái tạo TBM white-box ablation trên 20 RNA hiệu chỉnh.
  \item Tái tạo đánh giá Geometry v2 và lấy mẫu lại theo RNA với hạt giống ngẫu nhiên 2025.
  \item Tái tạo tập 60/20/20, xác nhận không có họ hoặc trình tự trùng giữa các tập.
  \item Chỉ huấn luyện trên tập huấn luyện; chỉ chọn mô hình và ngưỡng trên tập hiệu chỉnh.
  \item Mở tập xác nhận mới nhất đúng một lần và lưu kết quả bất kể chiều hướng.
  \item Kiểm tra định dạng tệp nộp cuối và nguồn gốc mọi dữ liệu đi kèm.
\end{enumerate}

\chapter{Các mục cần hoàn thiện trước khi nộp}

\begin{itemize}
  \item Thay toàn bộ marker \texttt{TODO} trên trang bìa và front matter.
  \item Chuẩn hóa mẫu LaTeX theo quy định của trường.
  \item Chốt tên luận văn tiếng Việt/Anh cùng giảng viên hướng dẫn.
  \item Chuyển tài liệu tham khảo nhúng sang BibTeX theo chuẩn của trường.
  \item Thêm hình mô tả dữ liệu, khung phương pháp và Geometry v2 từ kho mã.
  \item Bổ sung bảng kết quả từng RNA vào phụ lục nếu hội đồng yêu cầu.
  \item Bổ sung phần cứng, thời gian chạy, mã phiên bản Git và mã kiểm tra dữ liệu.
  \item Lưu mã phiên bản kernel, ID submission, manifest và mã kiểm tra SHA-256 của cả
  hai lần nộp Kaggle; không điều chỉnh ngược quy trình phát triển để khớp điểm tập riêng.
  \item Rà soát thuật ngữ Việt--Anh và định dạng số theo hướng dẫn khoa.
\end{itemize}

% -----------------------------------------------------------------------------
% REFERENCES — embedded to keep this as one self-contained .tex file.
% -----------------------------------------------------------------------------
\begin{thebibliography}{99}

\bibitem{tm_score}
Y. Zhang và J. Skolnick,
``Scoring function for automated assessment of protein structure template quality,''
\emph{Proteins}, 2004.

\bibitem{usalign}
J. Zhang và cộng sự,
``US-align: universal structure alignments of proteins, nucleic acids, and macromolecular complexes,''
\emph{Nature Methods}, 2022.

\bibitem{lddt}
V. Mariani, M. Biasini, A. Barbato và T. Schwede,
``lDDT: a local superposition-free score for comparing protein structures and models using distance difference tests,''
\emph{Bioinformatics}, 2013.

\bibitem{rna_puzzles}
J. A. Cruz và cộng sự,
``RNA-Puzzles: a CASP-like evaluation of RNA three-dimensional structure prediction,''
\emph{RNA}, 2012.

\bibitem{rna_puzzles_toolkit}
M. Magnus, M. Antczak, T. Zok và cộng sự,
``RNA-Puzzles toolkit: a computational resource of RNA 3D structure benchmark datasets,
structure manipulation, and evaluation tools,''
\emph{Nucleic Acids Research}, tập 48, số 2, trang 576--588, 2020,
doi: 10.1093/nar/gkz1108.

\bibitem{rna_progress_review}
X. Wang, S. Yu, E. Lou, Y.-L. Tan và Z.-J. Tan,
``RNA 3D Structure Prediction: Progress and Perspective,''
\emph{Molecules}, tập 28, số 14, bài 5532, 2023,
doi: 10.3390/molecules28145532.

\bibitem{state_of_rnart}
C. Bernard, G. Postic, S. Ghannay và F. Tahi,
``State-of-the-RNArt: benchmarking current methods for RNA 3D structure prediction,''
\emph{NAR Genomics and Bioinformatics}, tập 6, số 2, lqae048, 2024,
doi: 10.1093/nargab/lqae048.

\bibitem{computational_rna_review}
Y. Kagaya, B. Liu và D. Kihara,
``Computational approaches for RNA structure prediction and design,''
\emph{Cell Reports Physical Science}, tập 7, số 2, bài 103097, 2026,
doi: 10.1016/j.xcrp.2026.103097.

\bibitem{rna_composer}
M. Biesiada, K. J. Purzycka, M. Szachniuk, J. Blazewicz và R. W. Adamiak,
``Automated RNA 3D Structure Prediction with RNAComposer,''
\emph{Methods in Molecular Biology}, tập 1490, trang 199--215, 2016,
doi: 10.1007/978-1-4939-6433-8\_13.

\bibitem{drfold}
Y. Li, C. Zhang, C. Feng, R. Pearce, P. L. Freddolino và Y. Zhang,
``Integrating end-to-end learning with deep geometrical potentials for ab initio RNA
structure prediction,''
\emph{Nature Communications}, tập 14, bài 5745, 2023.

\bibitem{rhofoldplus}
T. Shen, Z. Hu, S. Sun và cộng sự,
``Accurate RNA 3D structure prediction using a language model-based deep learning approach,''
\emph{Nature Methods}, tập 21, trang 2287--2298, 2024,
doi: 10.1038/s41592-024-02487-0.

\bibitem{drfold2}
Y. Li, C. Feng, X. Zhang, S. Tsukiyama, D. Feng và Y. Zhang,
``DRfold2 is a deep learning-based tool that enables efficient and accurate RNA
structure prediction,''
\emph{PLOS Biology}, tập 24, số 2, e3003659, 2026,
doi: 10.1371/journal.pbio.3003659.

\bibitem{boltz1}
J. Wohlwend, G. Corso, S. Passaro và cộng sự,
``Boltz-1: Democratizing Biomolecular Interaction Modeling,''
\emph{bioRxiv}, 2024,
doi: 10.1101/2024.11.19.624167.

\bibitem{ares}
R. J. L. Townshend, S. Eismann, A. M. Watkins và cộng sự,
``Geometric deep learning of RNA structure,''
\emph{Science}, tập 373, trang 1047--1051, 2021,
doi: 10.1126/science.abe5650.

\bibitem{rnassess}
P. Lukasiak, M. Antczak, T. Ratajczak và J. Blazewicz,
``RNAssess---a web server for quality assessment of RNA 3D structures,''
\emph{Nucleic Acids Research}, tập 43, trang W502--W506, 2015,
doi: 10.1093/nar/gkv557.

\bibitem{kaggle_competition}
Stanford RNA 3D Folding Competition,
Kaggle competition documentation,
\url{https://www.kaggle.com/competitions/stanford-rna-3d-folding}.

\bibitem{top1_tbm}
JaeJohn,
``RNA 3D folds --- TBM only approach,'' Kaggle notebook, version được báo cáo 0.59298,
\url{https://www.kaggle.com/code/jaejohn/rna-3d-folds-tbm-only-approach}.

\bibitem{top1_hybrid}
JaeJohn,
``Sub 1 4 4 hybrid final take,'' Kaggle notebook, version được báo cáo 0.57631,
\url{https://www.kaggle.com/code/jaejohn/sub-1-4-4-hybrid-final-take}.

\end{thebibliography}

\end{document}
